% Modular main.tex: includes front-matter files from FrontBackMatter/
\documentclass[11pt,a4paper,oneside]{article}

\usepackage[utf8]{inputenc}
\usepackage[T1,T2A]{fontenc}
\usepackage[mongolian,english]{babel}

% Times New Roman фонт (бие)
\usepackage{times}

% ============================================================
% Geometry - ТСА-ын албан шаардлагын дагуу
% Зүүн тал: 3.0 cm, Баруун тал: 1.5 cm
% Дээд тал: 2.0 cm, Доод тал: 2.0 cm
% ============================================================
\usepackage{geometry}
\geometry{
  a4paper,
  left=3.0cm,
  right=1.5cm,
  top=2.0cm,
  bottom=2.0cm,
  headheight=14pt,
  headsep=10pt,
  footskip=25pt,
  includeheadfoot=false
}

% Мөр хоорондын зай 1.15, догол мөр өмнөх зай 6pt
\usepackage{setspace}
\setstretch{1.15}
\setlength{\parskip}{6pt}
\setlength{\parindent}{0pt}

% Зүүн/баруун талын текст зэрэгцүүлэлт (justify)
% Энэ нь текстийг хоёр талаас нь тэгш зэрэгцүүлж, хэт том зай гарахаас сэргийлнэ
\tolerance=1000
\emergencystretch=3em
\hyphenpenalty=50
\hbadness=10000

% Bibliography (biblatex + biber)
\usepackage[backend=biber,style=ieee,sorting=none]{biblatex}
\addbibresource{references.bib}

% Graphics, math, tables
\usepackage{graphicx,amsmath,amssymb,booktabs,caption,subcaption}
\usepackage{array,multirow,longtable}
\usepackage{multicol}
\usepackage{xcolor,colortbl}
\usepackage{float}

% Code listings
\usepackage{listings}
\definecolor{codegreen}{rgb}{0,0.6,0}
\definecolor{codegray}{rgb}{0.5,0.5,0.5}
\definecolor{codepurple}{rgb}{0.58,0,0.82}
\definecolor{backcolour}{rgb}{0.95,0.95,0.92}

\lstdefinelanguage{JavaScript}{
  keywords={typeof, new, true, false, catch, function, return, null, catch, switch, var, if, in, while, do, else, case, break, const, let, async, await, throw, import, export, from, default},
  keywordstyle=\color{blue}\bfseries,
  ndkeywords={class, export, boolean, throw, implements, import, this},
  ndkeywordstyle=\color{darkgray}\bfseries,
  identifierstyle=\color{black},
  sensitive=false,
  comment=[l]{//},
  morecomment=[s]{/*}{*/},
  commentstyle=\color{codegreen}\ttfamily,
  stringstyle=\color{codepurple}\ttfamily,
  morestring=[b]',
  morestring=[b]"
}

\lstdefinestyle{mystyle}{
    backgroundcolor=\color{backcolour},   
    commentstyle=\color{codegreen},
    keywordstyle=\color{magenta},
    numberstyle=\tiny\color{codegray},
    stringstyle=\color{codepurple},
    basicstyle=\ttfamily\footnotesize,
    breakatwhitespace=false,         
    breaklines=true,                 
    captionpos=b,                    
    keepspaces=true,                 
    numbers=left,                    
    numbersep=5pt,                  
    showspaces=false,                
    showstringspaces=false,
    showtabs=false,                  
    tabsize=2
}
\lstset{style=mystyle}

% TikZ for diagrams
\usepackage{tikz}
\usetikzlibrary{shapes,shapes.multipart,arrows,arrows.meta,positioning,fit,calc,shadows,backgrounds}

% PGFPlots for bar/line charts
\usepackage{pgfplots}
\pgfplotsset{compat=1.18}

% Links, TOC, header/footer
\usepackage[
    colorlinks=true,
    linkcolor=black,
    citecolor=black,
    filecolor=black,
    urlcolor=black,
    pdfborder={0 0 0}
]{hyperref}
\usepackage[titles]{tocloft}
\usepackage{fancyhdr}

% ToC Customization
\renewcommand{\cftsecleader}{\cftdotfill{\cftdotsep}}
\setcounter{tocdepth}{2}

% Top-level (section = бүлэг) оруулгуудыг bold болгох,
% subsection-ууд энгийн жинтэй үлдэх
% (cftsecfont-ыг PTSerif bold-оор дор дахин тодорхойлсон — энд subsec-ийг энгийн болгож байна)
\renewcommand{\cftsubsecfont}{\normalfont}
\renewcommand{\cftsubsecpagefont}{\normalfont}

\renewcommand{\cftaftertoctitle}{%
  \hfill\textbf{Хуудас}\par\medskip}

\addto\captionsenglish{\renewcommand{\listfigurename}{Зургийн жагсаалт}}
\addto\captionsenglish{\renewcommand{\listtablename}{Хүснэгтийн жагсаалт}}
\addto\captionsenglish{\renewcommand{\contentsname}{Гарчиг}}

\captionsetup[figure]{name=Зураг, font=small, labelfont=bf, justification=centering}
\captionsetup[table]{name=Хүснэгт, font=small, labelfont=bf, justification=centering}

\renewcommand{\cftfigleader}{\cftdotfill{\cftdotsep}}
\renewcommand{\cfttableader}{\cftdotfill{\cftdotsep}}

\renewcommand{\cftfigpresnum}{Зураг~}
\renewcommand{\cftfigaftersnum}{:}
\setlength{\cftfignumwidth}{4.5em}

\renewcommand{\cfttabpresnum}{Хүснэгт~}
\renewcommand{\cfttabaftersnum}{:}
\setlength{\cfttabnumwidth}{5.5em}

\renewcommand{\cftafterloftitle}{%
  \hfill\textbf{Хуудас}\par\medskip}
\renewcommand{\cftafterlottitle}{%
  \hfill\textbf{Хуудас}\par\medskip}

\usepackage{verbatim}

% Section / Subsection / Subsubsection гарчгуудыг bold болгох.
% Times нь T2A кириллийн bold хувилбар дэмждэггүй тул зөвхөн гарчгуудад
% PT Serif bold-ыг түр сэлгэж хэрэглэнэ. Биеийн текст Times хэвээр.
\usepackage{titlesec}
\newcommand{\headingfont}{\fontfamily{PTSerif-TLF}\selectfont\bfseries}
\titleformat{\section}{\headingfont\Large}{\thesection}{1em}{}
\titleformat{\subsection}{\headingfont\large}{\thesubsection}{1em}{}
\titleformat{\subsubsection}{\headingfont\normalsize}{\thesubsubsection}{1em}{}

% TOC дотор ч адил зөвхөн bold хэрэгтэй газарт PTSerif ашиглана
\renewcommand{\cftsecfont}{\fontfamily{PTSerif-TLF}\selectfont\bfseries}
\renewcommand{\cftsecpagefont}{\fontfamily{PTSerif-TLF}\selectfont\bfseries}

% Preamble metadata
\newcommand{\univname}{Шинэ Монгол Технологийн Коллеж}
\newcommand{\deptname}{Компьютерын ухааны тэнхим}
\newcommand{\ttitle}{Байгууллагын бүтээгдэхүүн хүргэлтийн нэгдсэн систем}
\newcommand{\ttitleeng}{Unified B2B Product Delivery System}
\newcommand{\thesistype}{ДЭД БАКАЛАВРЫН ТӨГСӨЛТИЙН СУДАЛГААНЫ АЖИЛ}

\newcommand{\authorshort}{Ү. Тамир}
\newcommand{\authorlong}{Үдэлгомбо Тамир}
\newcommand{\authorcode}{s21c011b}
\newcommand{\authoremail}{s21c011b@nmct.edu.mn}
\newcommand{\authorphone}{89475188}

\newcommand{\supervisor}{Н.Соронзонболд}

\newcommand{\submitdate}{2026}

\newcommand{\keywordsMN}{B2B платформ, захиалгын удирдлага, бодит цагийн хяналт, React Native, Supabase, Next.js, ePOD, гар утасны аппликейшн}
\newcommand{\keywordnames}{B2B платформ, захиалгын удирдлага, бодит цагийн хяналт, React Native, Supabase, Next.js, ePOD, гар утасны аппликейшн}
\newcommand{\authorshipname}{Зохиогчийн мэдүүлэг}
\newcommand{\acknowledgementname}{Талархал}
\newcommand{\abbrevname}{Товчлолын жагсаалт}
\newcommand{\addressname}{}

\newcommand{\shortname}{\authorshort}
\newcommand{\supname}{\supervisor}
\newcommand{\chairname}{\deptchair}
\newcommand{\readname}{\reader}

\newcommand{\addchaptertocentry}[1]{\addcontentsline{toc}{section}{#1}}
\newenvironment{declaration}{\begin{titlepage}}{\end{titlepage}}

% Header/Footer - текстийн өргөнтэй яг адил
\pagestyle{fancy}
\fancyhf{}
\fancyhead[L]{\footnotesize \ttitle}
\fancyhead[R]{\footnotesize \authorlong}
\fancyfoot[C]{\thepage}
\renewcommand{\headrulewidth}{0.5pt}
\renewcommand{\footrulewidth}{0.5pt}

% Document
\begin{document}

\pagenumbering{roman}

%--- FrontBackMatter/TitlePage.tex ---
\begin{titlepage}
\begin{center}
\vspace{-45mm}
\textbf{\Large\MakeUppercase{\univname}}\\[1.2mm]
\textbf{\Large\MakeUppercase{\deptname}}\\

\vspace{75mm}

\normalsize Дэд бакалаврын төгсөлтийн судалгааны ажил (061301)\\[2.5mm]
{\linespread{1.08}\selectfont\textbf{\Large \parbox{0.95\textwidth}{\centering \MakeUppercase{\ttitle}}}\\}


\vspace{20mm}
\begin{center}
\begin{minipage}{0.4\textwidth}
Удирдагч багш\\[1.5mm]
Гүйцэтгэсэн оюутан
\end{minipage}
\hspace{0.5cm}
\begin{minipage}{0.3\textwidth}
\raggedleft
/\supervisor, Магистр/\\[1.5mm]
/\authorcode, \authorshort/
\end{minipage}
\end{center}

\end{center}

\vspace{60mm}

\vfill
\begin{center}
 Улаанбаатар хот \hspace{1pt} \\
 \submitdate \hspace{1pt} он\\
\end{center}

\end{titlepage}

% FrontBackMatter/Abstract.tex
\addchaptertocentry{Хураангуй}

\vspace*{-1cm}
\vspace{0.8cm}

\begin{center}
{\headingfont\Large \ttitle\par}
\bigskip
{\headingfont\Large \deptname\par}
\bigskip

{\large \shortname\par}
{\headingfont\Large Хураангуй\par}

{\normalsize \par}
\end{center}
\vspace{1cm}


\bigskip

Монгол Улсын жижиг, дунд бизнесүүд хоорондоо бараа захиалах, хүргэх үйлээ гол төлөв утас, мессеж, цаасаар гүйцэтгэсээр байна. Ийм гараар зохицуулдаг хэлбэр нь цаг алдагдуулж, алдаа үүсгэж, явцыг хянахад төвөгтэй болгож байна.

Төсөлд зуучлагч оператор шаардахгүй, олон байгууллага зэрэг үйлчлүүлдэг, хүргэгчид өөрсдөө даалгавраа сонгодог B2B хүргэлтийн нэгдсэн платформыг боловсруулсан. Систем нь хүргэгчдэд зориулсан гар утасны аппликейшн, борлуулагчдад зориулсан веб аппликейшн, үүлэн дээр ажилладаг backend гэсэн гурван үндсэн хэсгээс бүрдэнэ.

Системийн найдвартай байдал, аюулгүй байдал, хурдны үзүүлэлтийг туршилтаар үнэлсэн. Нэг даалгаврыг олон хүргэгч зэрэг авахаар өрсөлдөхөд давхардал үүсэхгүй, байгууллагуудын мэдээлэл бие биендээ алдагдахгүй, хүргэлтийг баталгаажуулах код нь аюулгүй ажиллаж байгаа нь нотлогдсон.

Судалгааны үр дүнд тус систем нь Монголын жижиг, дунд бизнесийн хүргэлтийн гол бэрхшээлүүдийг бодитоор шийдэж чадах нь батлагдлаа.


\vfill

%-------------------------------------------------------------------------------
%	ABBREVIATIONS
%-------------------------------------------------------------------------------

% (Энийг main.tex дээр нэг удаа зарласан бол эндээс устгаж болно)
% \newcommand{\abbrevname}{Товчилсон үгийн жагсаалт}

\section*{\abbrevname}
\addcontentsline{toc}{section}{\abbrevname}

\vspace{1cm}

\begin{center}
\begin{tabular}{ll}
\textbf{ЖДБ} & Жижиг, дунд бизнес \\
\textbf{B2B} & \textbf{B}usiness \textbf{t}o \textbf{B}usiness \\
\textbf{API} & \textbf{A}pplication \textbf{P}rogramming \textbf{I}nterface \\
\textbf{BaaS} & \textbf{B}ackend-\textbf{a}s-\textbf{a}-\textbf{S}ervice \\
\textbf{BFF} & \textbf{B}ackend-\textbf{f}or-\textbf{F}rontend \\
\textbf{DB} & \textbf{D}ata\textbf{b}ase \\
\textbf{ERD} & \textbf{E}ntity \textbf{R}elationship \textbf{D}iagram \\
\textbf{RLS} & \textbf{R}ow \textbf{L}evel \textbf{S}ecurity \\
\textbf{RBAC} & \textbf{R}ole-\textbf{B}ased \textbf{A}ccess \textbf{C}ontrol \\
\textbf{RPC} & \textbf{R}emote \textbf{P}rocedure \textbf{C}all \\
\textbf{CRUD} & \textbf{C}reate, \textbf{R}ead, \textbf{U}pdate, \textbf{D}elete \\
\textbf{SQL} & \textbf{S}tructured \textbf{Q}uery \textbf{L}anguage \\
\textbf{JWT} & \textbf{J}SON \textbf{W}eb \textbf{T}oken \\
\textbf{UUID} & \textbf{U}niversally \textbf{U}nique \textbf{I}dentifier \\
\textbf{POD} & \textbf{P}roof \textbf{o}f \textbf{D}elivery \\
\textbf{ePOD} & \textbf{e}lectronic \textbf{P}roof \textbf{o}f \textbf{D}elivery \\
\textbf{OTP} & \textbf{O}ne-\textbf{T}ime \textbf{P}assword \\
\textbf{KYC} & \textbf{K}now \textbf{Y}our \textbf{C}ustomer \\
\textbf{FCFS} & \textbf{F}irst-\textbf{C}ome, \textbf{F}irst-\textbf{S}erved \\
\textbf{MVCC} & \textbf{M}ulti-\textbf{V}ersion \textbf{C}oncurrency \textbf{C}ontrol \\
\textbf{WAL} & \textbf{W}rite-\textbf{A}head \textbf{L}og \\
\textbf{SDK} & \textbf{S}oftware \textbf{D}evelopment \textbf{K}it \\
\textbf{SLA} & \textbf{S}ervice \textbf{L}evel \textbf{A}greement \\
\textbf{MVP} & \textbf{M}inimum \textbf{V}iable \textbf{P}roduct \\
\textbf{SSR} & \textbf{S}erver-\textbf{S}ide \textbf{R}endering \\
\textbf{ISR} & \textbf{I}ncremental \textbf{S}tatic \textbf{R}egeneration \\
\textbf{VU} & \textbf{V}irtual \textbf{U}ser (k6 load testing) \\
\textbf{RTT} & \textbf{R}ound-\textbf{T}rip \textbf{T}ime \\
\textbf{E2E} & \textbf{E}nd-\textbf{t}o-\textbf{E}nd \\
\textbf{CI} & \textbf{C}ontinuous \textbf{I}ntegration \\
\end{tabular}
\end{center}

\vfill
\clearpage
% =============================================================================
%   АЖЛЫН ТӨЛӨВЛӨГӨӨ
%   Төсөл: Байгууллагын бүтээгдэхүүн хүргэлтийн нэгдсэн систем
%   Оюутан: Ү.Тамир (s21c011b)
%   Удирдагч: Н.Соронзонболд
% =============================================================================

\addchaptertocentry{Ажлын төлөвлөгөө}

\begin{center}
{\Large\textbf{Ажлын төлөвлөгөө}}
\end{center}

\vspace{0.3cm}

\noindent\textbf{Оюутны нэр:} \authorshort \hfill \textbf{Удирдагч багшийн нэр:} \supervisor

\noindent\textbf{Судалгааны ажлын сэдэв:} \ttitle

\vspace{0.5cm}

\begin{table}[h!t]
\centering
\small
\begin{tabular}{|c|c|p{7.5cm}|c|}
\hline
\textbf{Сар} & \textbf{Долоо хоног} & \textbf{Төлөвлөгөө} & \textbf{Гүйцэтгэл} \\
\hline

% ----------------------------------------------------------
% 10 САР: Сэдэв сонголт, судалгаа, шинжилгээ
% ----------------------------------------------------------
\multirow{4}{*}{10 сар}
& 1 & Сэдэв сонгох, B2B хүргэлтийн зах зээлийн судалгаа эхлэх & $\checkmark$ \\
\cline{2-4}
& 2 & Last-mile логистик, crowdsourced хүргэлтийн онолын судалгаа & $\checkmark$ \\
\cline{2-4}
& 3 & Олон улсын платформ судалгаа (Uber Freight) & $\checkmark$ \\
\cline{2-4}
& 4 & Дотоодын платформ судалгаа (Papa Logistics), харьцуулалт & $\checkmark$ \\
\hline

% ----------------------------------------------------------
% 11 САР: Архитектурын загварчлал, ӨС дизайн
% ----------------------------------------------------------
\multirow{4}{*}{11 сар}
& 1 & Системийн архитектур төлөвлөх, технологийн стек сонгох & $\checkmark$ \\
\cline{2-4}
& 2 & Өгөгдлийн сангийн ERD загвар боловсруулах (11 хүснэгт) & $\checkmark$ \\
\cline{2-4}
& 3 & Supabase төсөл үүсгэх, PostgreSQL хүснэгтүүд хөгжүүлэх & $\checkmark$ \\
\cline{2-4}
& 4 & RLS бодлого, триггер, серверийн функцүүд хэрэгжүүлэх & $\checkmark$ \\
\hline

% ----------------------------------------------------------
% 12 САР: ТСА Үзлэг 1 + Backend бүрэн хөгжүүлэлт
% ----------------------------------------------------------
\multirow{4}{*}{12 сар}
& 1 & \textcolor{red}{ТСА Үзлэг 1} & $\checkmark$ \\
\cline{2-4}
& 2 & Supabase Auth, Storage тохиргоо, KYC баталгаажуулалтын процесс & $\checkmark$ \\
\cline{2-4}
& 3 & Хүргэгчийн мобайл апп (React Native/Expo) суурь бүтэц & $\checkmark$ \\
\cline{2-4}
& 4 & Бүртгэл, нэвтрэлт, профайл удирдлагын дэлгэцүүд & $\checkmark$ \\
\hline

% ----------------------------------------------------------
% 1 САР: Мобайл апп + Веб апп хөгжүүлэлт
% ----------------------------------------------------------
\multirow{4}{*}{1 сар}
& 1 & Хүргэгчийн апп: нээлттэй даалгавар, даалгавар авах функц & $\checkmark$ \\
\cline{2-4}
& 2 & Хүргэгчийн апп: хүргэлт гүйцэтгэх, орлого хянах & $\checkmark$ \\
\cline{2-4}
& 3 & Борлуулагчийн веб апп (Next.js) суурь бүтэц, бүтээгдэхүүн удирдлага & $\checkmark$ \\
\cline{2-4}
& 4 & Веб апп: захиалга удирдлага, хүргэлтийн даалгавар үүсгэх & $\checkmark$ \\
\hline

% ----------------------------------------------------------
% 2 САР: Туршилт + Тайлан бичиж эхлэх
% ----------------------------------------------------------
\multirow{4}{*}{2 сар}
& 1 & Бодит цагийн мэдэгдэл, нэгтгэлийн туршилт & $\checkmark$ \\
\cline{2-4}
& 2 & Аюулгүй байдлын тест, зэрэгцээ хандалтын тест & $\checkmark$ \\
\cline{2-4}
& 3 & Тайлангийн 1-р бүлэг (Удиртгал) бичих & $\checkmark$ \\
\cline{2-4}
& 4 & Тайлангийн 2-р бүлэг (Судалгаа, онол) бичих & $\checkmark$ \\
\hline

% ----------------------------------------------------------
% 3 САР: ТСА Үзлэг 2 + Тайлан үргэлжлүүлэх
% ----------------------------------------------------------
\multirow{4}{*}{3 сар}
& 1 & \textcolor{red}{ТСА Үзлэг 2} & $\checkmark$ \\
\cline{2-4}
& 2 & Тайлангийн 3-р бүлэг (Арга зүй, диаграмм) бичих & $\checkmark$ \\
\cline{2-4}
& 3 & Тайлангийн 4-р бүлэг (Дүгнэлт) бичих & $\checkmark$ \\
\cline{2-4}
& 4 & Санал хүсэлтийн дагуу засвар хийх, код баримтжуулалт & $\checkmark$ \\
\hline

% ----------------------------------------------------------
% 4 САР: Тайлан эцэслэх + Урьдчилсан хамгаалалтын бэлтгэл
% ----------------------------------------------------------
\multirow{4}{*}{4 сар}
& 1 & График, хүснэгт, диаграммуудыг эцэслэх & $\checkmark$ \\
\cline{2-4}
& 2 & Ном зүй, хавсралт, товчлол бэлдэх & $\checkmark$ \\
\cline{2-4}
& 3 & Хамгаалалтын слайд бэлдэх, дадлага хийх & $\checkmark$ \\
\cline{2-4}
& 4 & \textcolor{red}{ТСА урьдчилсан хамгаалалт} & $\square$ \\
\hline

% ----------------------------------------------------------
% 5 САР: Эцсийн засвар + Үндсэн хамгаалалт
% ----------------------------------------------------------
\multirow{4}{*}{5 сар}
& 1 & \textcolor{red}{ТСА урьдчилсан хамгаалалт} & $\square$ \\
\cline{2-4}
& 2 & Урьдчилсан хамгаалалтын санал тусгах, засвар хийх & $\square$ \\
\cline{2-4}
& 3 & Слайд сайжруулах, эцсийн бэлтгэл & $\square$ \\
\cline{2-4}
& 4 & \textcolor{red}{ТСА үндсэн хамгаалалт} & $\square$ \\
\hline

\end{tabular}
\end{table}

\vspace{2mm}
\noindent{\small Гүйцэтгэлийн төлөв шинэчлэгдсэн огноо: 2026-04-19.}

\vspace{8mm}
\noindent\begin{minipage}{0.4\textwidth}
Удирдагч багш:
\end{minipage}
\hfill
\begin{minipage}{0.4\textwidth}
/\supervisor/
\end{minipage}

\vspace{5mm}
\noindent\begin{minipage}{0.4\textwidth}
Гүйцэтгэсэн оюутан:
\end{minipage}
\hfill
\begin{minipage}{0.4\textwidth}
/\authorcode, \authorshort/
\end{minipage}
%-------------------------------------------------------------------------------
%	TABLE OF CONTENTS
%-------------------------------------------------------------------------------

\renewcommand{\contentsname}{Гарчиг}

\phantomsection
\addcontentsline{toc}{section}{Гарчиг}
\tableofcontents

\vfill
\clearpage

%-------------------------------------------------------------------------------
%	ЗУРГИЙН ЖАГСААЛТ
%-------------------------------------------------------------------------------
\phantomsection
\addcontentsline{toc}{section}{Зургийн жагсаалт}
\listoffigures
\clearpage

%-------------------------------------------------------------------------------
%	ХҮСНЭГТИЙН ЖАГСААЛТ
%-------------------------------------------------------------------------------
\phantomsection
\addcontentsline{toc}{section}{Хүснэгтийн жагсаалт}
\listoftables
\clearpage

\cleardoublepage

\pagenumbering{arabic}
\setcounter{page}{1}

\newpage

%-------------------------------------------------------------------------------
%   CHAPTER 1 - INTRODUCTION
%-------------------------------------------------------------------------------

\section{Удиртгал}
\label{sec:introduction}

\subsection{Асуудлын тулгамдсан байдал}

Монгол Улсын жижиг, дунд бизнесийн байгууллагууд (ЖДБ) хоорондоо бараа бүтээгдэхүүн захиалах, хүргэх үйл явцаа өдгөө хүртэл гол төлөв утас, мессеж, цаасан бичгээр гардаг~\cite{nso2023}. Захиалагч барааг өөрийн биеэр авч ирэх, эсхүл хүргэлтийг тусад нь зохион байгуулах шаардлага нь цаг алдагдуулж, зохион байгуулалтын алдаа үүсгэж, үйл явцыг хянахад төвөгтэй болгож байна.

Дижитал шилжилтийн нөхцөлд B2B харилцааг автоматжуулах эрэлт дэлхий нийтэд эрс нэмэгдэж~\cite{mckinsey2021delivery, grandview2023}, хүргэлтийн цахим платформууд ($\approx$4.5 тэрбум ам.доллар, 2023 он) гол хөдөлгөгч хүч болж байна~\cite{wang2016bigdata}. Монголын зах зээлд энэ чиглэлийн нэгдсэн шийдэл хараахан төлөвшөөгүй бөгөөд орон нутгийн тоглогчид (Papa Logistics гэх мэт) голдуу dispatcher-т тулгуурласан, нэг талтай загвартай хэвээр байна.

\subsection{Судалгааны асуулт}

Энэ ажил нь \textbf{operator-гүй, multi-tenant B2B хүргэлтийн платформыг одоогийн cloud BaaS технологи дээр найдвартай хэрхэн бүтээх вэ?} гэсэн ерөнхий асуултыг дэвшүүлж, гурван тодорхой дэд асуултыг шалгана:

\begin{itemize}
    \item \textbf{RQ1 — Concurrency:} PostgreSQL-ийн \texttt{FOR UPDATE SKIP LOCKED} механизм нь FCFS claim model-ыг race condition-гүй хэрэгжүүлж чадах уу?
    \item \textbf{RQ2 — Isolation:} Row Level Security (RLS) нь application-layer алдаанаас үл хамааран multi-tenant өгөгдлийн тусгаарлалтыг database түвшинд хангаж чадах уу?
    \item \textbf{RQ3 — Viability:} BaaS-mediated архитектур нь зөвхөн prototype биш, production-д хүрэх системийг хурдан хөгжүүлэх суурь болж чадах уу?
\end{itemize}

\subsection{Зорилго, зорилт}

Ажлын \textbf{зорилго} нь B2B хүргэлтийн operator-гүй, multi-tenant платформыг Supabase (PostgreSQL + PostgREST + RLS) технологи дээр боловсруулж, гурван дэд асуултыг туршилтаар хариулах явдал юм.

Дэвшүүлсэн \textbf{зорилтууд}:

\begin{enumerate}
    \item \textbf{Судалгаа, шинжилгээ:} Монгол Улс дахь B2B хүргэлтийн өнөөгийн байдлыг дүгнэх, Uber Freight, Papa Logistics зэрэг ижил төрлийн платформуудыг харьцуулах, архитектурын зөрүү ба design шаардлагыг тодорхойлох.
    \item \textbf{Архитектур, загварчлал:} Client layer (React Native/Expo + Next.js), BaaS layer (Supabase), data layer (PostgreSQL + RLS)-ээс тогтсон BaaS-mediated архитектур зохион бүтээх. 11 хүснэгт бүхий schema, ERD, sequence diagram, use case diagram боловсруулах.
    \item \textbf{Аюулгүй байдал:} RBAC, RLS, JWT authentication, bcrypt ePOD OTP-г нэгтгэсэн олон давхаргат security хэрэгжүүлэх.
    \item \textbf{Функционал модулиуд:} Захиалга, хүргэлтийн даалгавар, marketplace, realtime мэдэгдэл, ePOD, орлого, санхүүгийн тайланг бүрэн хэрэгжүүлэх.
\end{enumerate}

\subsection{Хамрах хүрээ ба хязгаарлалт}

Ажлын хамрах хүрээ нь \textbf{(i)} Монгол Улсын B2B хүргэлтийн orchestration, \textbf{(ii)} дунд зэргийн ачаалал (100 concurrent user хүртэл) бүхий staging орчинд хэрэгжсэн техникийн нотолгоо, \textbf{(iii)} single-region Supabase Pro tier тохиргоогоор хязгаарлагдана. Бодит GPS location tracking, төлбөрийн бодит integration (QPay, SocialPay), ML recommendation нь энэ ажлын хүрээнд биш — цаашдын хөгжүүлэлтийн чиглэл юм (§\ref{subsec:results_summary}).

\subsection{Ажлын бүтэц}

Үлдэж буй хэсэг нь дараах бүтэцтэй. \textbf{Бүлэг 2}-т last-mile логистик, two-sided marketplace, multi-tenant архитектур, concurrency control, BaaS-ын онолын суурь болон Uber Freight, Papa Logistics-тэй харьцуулсан судалгааг танилцуулна. \textbf{Бүлэг 3}-т системийн архитектур, өгөгдлийн загвар, аюулгүй байдлын бодлого, функционал шаардлагыг дэлгэрэнгүй тайлбарлана. \textbf{Бүлэг 4}-т гүйцэтгэсэн туршилтын үр дүн, харьцуулалт болон дүгнэлтийг оруулсан.

%-------------------------------------------------------------------------------
%   CHAPTER 2 - RESEARCH
%-------------------------------------------------------------------------------

\section{Онолын болон харьцуулсан судалгаа}

\subsection{Онолын судалгаа}
Энэ хэсэгт B2B түгээлтийн бизнесийн процесс, логистик систем, маршрутын оновчлол, болон гар утасны хөгжүүлэлтийн технологиудын талаарх онолын ойлголтуудыг дэлгэрэнгүй авч үзнэ.

\subsubsection{Нийлүүлэлтийн сүлжээний удирдлага (SCM)}
Нийлүүлэлтийн сүлжээний удирдлага (Supply Chain Management - SCM) нь түүхий эд бэлтгэн нийлүүлэгчээс эхлээд эцсийн хэрэглэгч хүртэлх бараа, үйлчилгээ, мэдээлэл, санхүүгийн урсгалыг төлөвлөх, хэрэгжүүлэх, хянах цогц үйл ажиллагаа юм. Орчин үед SCM нь зөвхөн ложистик бус, мэдээллийн технологитой нягт уялдаж, "Ухаалаг нийлүүлэлтийн сүлжээ" (Smart Supply Chain) болон хөгжиж байна. SCM-ийн үндсэн зорилго нь нийлүүлэлтийн сүлжээний нийт зардлыг бууруулахын зэрэгцээ хэрэглэгчийн сэтгэл ханамжийг дээд зэргээр хангахад оршино.

\textbf{SCM-ийн үндсэн бүрэлдэхүүн хэсгүүд:}
\begin{itemize}
    \item \textbf{Төлөвлөлт (Planning):} Эрэлтийг таамаглах, нөөцийг төлөвлөх.
    \item \textbf{Худалдан авалт (Sourcing):} Нийлүүлэгчийг сонгох, гэрээ байгуулах.
    \item \textbf{Үйлдвэрлэл (Making):} Түүхий эдийг бэлэн бүтээгдэхүүн болгох.
    \item \textbf{Түгээлт (Delivering):} Логистик, агуулах, тээвэрлэлт.
    \item \textbf{Буцаалт (Returning):} Гэмтэлтэй эсвэл илүүдэл барааг буцаах үйл явц.
\end{itemize}

\subsubsection{B2B Түгээлтийн процесс ба онцлог}
B2B (Business-to-Business) түгээлт нь үйлдвэрлэгчээс бөөний худалдаачид, эсвэл бөөний худалдаачдаас жижиглэн худалдаачид руу бараа нийлүүлэх үйл явц юм. B2C (Business-to-Consumer) хүргэлтээс ялгаатай нь:
\begin{itemize}
    \item \textbf{Захиалгын хэмжээ:} B2B захиалга нь ихэвчлэн тоо хэмжээний хувьд их, овор хэмжээ томтой байдаг.
    \item \textbf{Давтамж:} Тогтмол хуваарийн дагуу, давтагдах шинж чанартай (жишээ нь: долоо хоног бүрийн Даваа гарагт).
    \item \textbf{Харилцаа:} Урт хугацааны гэрээт харилцаанд суурилдаг бөгөөд итгэлцэл чухал үүрэгтэй.
    \item \textbf{Шийдвэр гаргалт:} Олон шат дамжлагатай, оновчтой байдлыг чухалчилдаг. Худалдан авалтын шийдвэрийг сэтгэл хөдлөлөөр бус, үр ашгийн тооцооллоор гаргадаг.
\end{itemize}

\subsubsection{Last-Mile Delivery (Сүүлийн цэгийн хүргэлт)}
"Last-Mile Delivery" гэдэг нь барааг түгээлтийн төвөөс эцсийн хэрэглэгч (дэлгүүр эсвэл хувь хүн) хүртэл хүргэх логистикийн сүүлийн үе шатыг хэлнэ. Энэ нь нийлүүлэлтийн сүлжээний хамгийн өндөр зардалтай (нийт зардлын 53\% хүртэл), хамгийн их цаг хугацаа шаардсан, хамгийн төвөгтэй хэсэгт тооцогддог. Улаанбаатар хотын нөхцөлд замын түгжрэл, хаягжилтын асуудал нь Last-Mile Delivery-ийн үр ашгийг эрс бууруулдаг.

\subsubsection{Маршрутын оновчлол (Route Optimization)}
Хүргэлтийн зардлыг бууруулах, цаг хугацааг хэмнэхэд маршрутын оновчлол чухал үүрэгтэй. Энэ нь "Traveling Salesman Problem" (TSP) буюу олон цэгт хүргэлт хийх хамгийн дөт, үр ашигтай замыг тооцоолох математик аргачлал дээр суурилдаг. Орчин үед Google Maps API, Mapbox зэрэг технологиудыг ашиглан замын түгжрэл, цаг агаар, тээврийн хэрэгслийн даац зэрэг олон хүчин зүйлийг тооцон маршрутыг динамикаар төлөвлөх боломжтой болсон.

\subsubsection{Электрон баримт бичиг (e-POD)}
Electronic Proof of Delivery (e-POD) буюу Цахим хүргэлтийн баталгаажуулалт нь цаасан баримтыг халж, дижитал хэлбэрээр хүргэлтийг баталгаажуулах стандарт юм. e-POD систем нь дараах мэдээллийг агуулдаг:
\begin{itemize}
    \item Хүлээн авагчийн цахим гарын үсэг.
    \item Хүргэгдсэн барааны зураг.
    \item Хүргэлт хийгдсэн газарзүйн байршил (GPS координат).
    \item Хүргэлт хийгдсэн цаг хугацаа (Timestamp).
\end{itemize}
Энэ нь маргаантай асуудлыг шийдвэрлэх, төлбөр тооцоог хурдасгахад чухал ач холбогдолтой.

\subsection{Технологийн судалгаа}
Энэхүү төслийг хэрэгжүүлэхэд ашигласан үндсэн технологиудын онцлог, давуу талуудыг судлав. Tdelivery системийн хөгжүүлэлтэд ашигласан технологиудыг дараах ангиллаар танилцуулав. Орчин үеийн гар утасны аппликейшн хөгжүүлэлтэд олон технологи хамтран ажилладаг бөгөөд тус бүрийн үүрэг, ач холбогдлыг дэлгэрэнгүй авч үзье.

\subsubsection{Програмчлалын хэл}
Системийг хөгжүүлэхэд дараах програмчлалын хэлүүдийг ашигласан. Програмчлалын хэл нь аппликейшний суурь болох бөгөөд зөв сонголт хийх нь төслийн амжилтад шийдвэрлэх нөлөөтэй.

\begin{table}[H]
\centering
\caption{Програмчлалын хэлүүд}
\label{tab:programming-languages}
\begin{tabular}{|l|l|p{7cm}|}
\hline
\textbf{Технологи} & \textbf{Хувилбар} & \textbf{Тайлбар} \\
\hline
TypeScript & \textasciitilde5.9.2 & Үндсэн програмчлалын хэл. JavaScript-ийн супер сет бөгөөд статик төрөл шалгалттай. \\
\hline
JavaScript (ES6+) & - & React Native-ийн суурь хэл \\
\hline
JSX/TSX & - & UI компонентуудыг тодорхойлох синтакс \\
\hline
CSS & - & Стиль (Tailwind CSS-ээр дамжуулан) \\
\hline
\end{tabular}
\end{table}

\textbf{TypeScript гэж юу вэ?}

TypeScript нь Microsoft компаниас 2012 онд хөгжүүлсэн програмчлалын хэл бөгөөд JavaScript-ийн "супер сет" (superset) юм. Энэ нь JavaScript-ийн бүх боломжийг агуулахын зэрэгцээ статик төрөл шалгалт (static type checking), интерфэйс (interfaces), энумераци (enums), женерик (generics) зэрэг нэмэлт боломжуудыг санал болгодог. TypeScript код нь компайл хийгдэж цэвэр JavaScript болдог тул аль ч JavaScript орчинд ажиллах боломжтой.

\textbf{TypeScript-ийн давуу талууд:}
\begin{itemize}
    \item \textbf{Статик төрөл шалгалт (Static Type Checking):} Компайл хийх үед алдааг илрүүлж, runtime алдааг багасгадаг. Жишээ нь: \texttt{let count: number = "hello"} гэж бичвэл компайл хийх үед алдаа гарна. Энэ нь runtime үед гарах алдааг урьдчилан сэргийлдэг.
    \item \textbf{Код дахин ашиглалт:} Interface болон Type-ууд нь кодын бүтцийг тодорхой болгож, дахин ашиглах боломжийг нэмэгдүүлдэг. Жишээ нь захиалгын бүтцийг нэг удаа тодорхойлоод бүх газар ашиглах боломжтой.
    \item \textbf{IDE дэмжлэг:} VS Code зэрэг орчин үеийн IDE-үүдэд автомат бөглөлт (IntelliSense), алдаа илрүүлэлт, рефакторинг зэрэг боломжуудыг сайн дэмждэг. Энэ нь хөгжүүлэлтийн бүтээмжийг 20-30\% нэмэгдүүлдэг гэж үздэг.
    \item \textbf{Том төслийн удирдлага:} Олон хөгжүүлэгчтэй төсөлд кодын чанарыг хадгалахад тустай. Төрлийн тодорхойлолт нь баримт бичиг шиг үүрэг гүйцэтгэдэг.
    \item \textbf{Алдааг эрт илрүүлэх:} JavaScript-д runtime үед гардаг алдаануудын 15\% орчимыг компайл хийх үед илрүүлдэг.
\end{itemize}

\textbf{JavaScript (ES6+) ба JSX/TSX}

JavaScript нь вэб болон гар утасны аппликейшн хөгжүүлэлтийн хамгийн түгээмэл хэл бөгөөд ES6 (ECMAScript 2015) хувилбараас хойш олон чухал шинэчлэлтүүд нэмэгдсэн:
\begin{itemize}
    \item \textbf{Arrow Functions:} Богино синтаксээр функц тодорхойлох (\texttt{const add = (a, b) => a + b})
    \item \textbf{Async/Await:} Асинхрон кодыг синхрон маягаар бичих боломж
    \item \textbf{Destructuring:} Объект болон массивын утгуудыг задлах (\texttt{const \{name, age\} = user})
    \item \textbf{Spread Operator:} Массив, объектыг тараах (\texttt{[...arr1, ...arr2]})
    \item \textbf{Template Literals:} Тэмдэгт мөрийг динамикаар үүсгэх (\texttt{`Hello \$\{name\}`})
\end{itemize}

JSX (JavaScript XML) нь React-д UI компонентуудыг HTML-тэй төстэй синтаксээр бичих боломжийг олгодог. TSX нь TypeScript дээр JSX ашиглах өргөтгөл юм. Жишээ нь:
\begin{lstlisting}[language=JavaScript]
// TSX жишээ
const OrderCard: React.FC<{order: Order}> = ({order}) => (
  <View style={styles.card}>
    <Text>{order.customerName}</Text>
    <Text>{order.totalAmount}</Text>
  </View>
);
\end{lstlisting}

\subsubsection{Framework ба Platform}
React Native болон Expo платформ нь гар утасны аппликейшн хөгжүүлэлтийн үндэс болсон. Эдгээр технологиуд нь нэг кодын баазаар олон платформд ажиллах боломжийг олгодог.

\begin{table}[H]
\centering
\caption{Framework ба Platform}
\label{tab:framework-platform}
\begin{tabular}{|l|l|p{7cm}|}
\hline
\textbf{Технологи} & \textbf{Хувилбар} & \textbf{Тайлбар} \\
\hline
React Native & 0.81.5 & Гар утасны апп хөгжүүлэлтийн фреймворк \\
\hline
React & 19.1.0 & UI компонент сан \\
\hline
Expo & \textasciitilde54.0.20 & React Native хөгжүүлэлтийн платформ \\
\hline
React Native Web & \textasciitilde0.21.0 & Веб дэмжлэг \\
\hline
\end{tabular}
\end{table}

\textbf{React Native гэж юу вэ?}

React Native нь Facebook (Meta) компаниас 2015 онд нээлттэй эх болгон гаргасан гар утасны аппликейшн хөгжүүлэлтийн фреймворк юм. React Native нь JavaScript ашиглан бичигдсэн кодыг native компонентууд руу хөрвүүлдэг тул бусад hybrid шийдлүүдээс (Cordova, Ionic) илүү өндөр гүйцэтгэлтэй. 

\textbf{React Native-ийн ажиллах зарчим:}
\begin{enumerate}
    \item JavaScript код нь JavaScript Engine (Hermes эсвэл JSC) дээр ажиллана.
    \item Bridge (гүүр) нь JavaScript болон Native код хооронд мэдээлэл дамжуулна.
    \item Native модулиуд нь төхөөрөмжийн боломжуудад (камер, GPS, гэх мэт) хандана.
    \item UI компонентууд нь жинхэнэ native элементүүд рүү хөрвүүлэгдэнэ.
\end{enumerate}

\textbf{React Native-ийн давуу талууд:}
\begin{itemize}
    \item \textbf{Cross-Platform (Олон платформ):} Нэг кодын баазаар iOS болон Android үйлдлийн системд зэрэг хөгжүүлэлт хийх боломжтой. Судалгаагаар кодын 80-90\%-ийг дахин ашиглах боломжтой. Энэ нь хөгжүүлэлтийн цаг болон зардлыг 40-50\% хэмнэдэг.
    \item \textbf{Performance (Гүйцэтгэл):} Native компонентүүдийг ашигладаг тул Hybrid аппликейшнүүдтэй (WebView суурьтай) харьцуулахад 60 FPS гүйцэтгэлтэй, анимаци гөлгөр ажилладаг.
    \item \textbf{Hot Reloading (Халуун дахин ачаалалт):} Кодод өөрчлөлт оруулахад аппликейшнийг дахин ачаалахгүйгээр үр дүнг шууд харах боломжтой. Энэ нь хөгжүүлэлтийн хурдыг 2-3 дахин нэмэгдүүлдэг.
    \item \textbf{Том коммьюнити:} Дэлхий даяар 2 сая гаруй хөгжүүлэгч ашигладаг тул тусламж, бэлэн сангууд (npm дээр 400,000+ package) ихтэй.
    \item \textbf{Код дахин ашиглалт:} Бизнес логик, утилит функцууд, төлөв удирдлага зэрэг кодыг iOS, Android, Web хооронд бүрэн дахин ашиглах боломжтой.
\end{itemize}

\textbf{React гэж юу вэ?}

React нь хэрэглэгчийн интерфэйс (UI) бүтээхэд зориулагдсан JavaScript сан бөгөөд Facebook-аас 2013 онд гаргасан. React-ийн гол онцлог нь компонент суурьтай архитектур юм. UI-г жижиг, дахин ашиглагдах компонентуудад хуваан зохион байгуулдаг.

\textbf{React-ийн гол ойлголтууд:}
\begin{itemize}
    \item \textbf{Virtual DOM:} Бодит DOM-ийг шууд өөрчлөхийн оронд виртуал DOM дээр өөрчлөлт хийгээд зөвхөн өөрчлөгдсөн хэсгийг бодит DOM-д шинэчилдэг. Энэ нь гүйцэтгэлийг эрс нэмэгдүүлдэг.
    \item \textbf{Component Lifecycle:} Компонент үүсэх, шинэчлэгдэх, устах үеийн амьдралын мөчлөгийн удирдлага.
    \item \textbf{Hooks:} \texttt{useState}, \texttt{useEffect}, \texttt{useContext} зэрэг функцууд нь функциональ компонентуудад төлөв болон хажуугийн нөлөө (side effects) удирдах боломжийг олгодог.
    \item \textbf{Props ба State:} Props нь эцэг компонентоос хүүхэд компонент руу дамжуулагдах өгөгдөл, State нь компонентын дотоод төлөв юм.
\end{itemize}

\textbf{Expo гэж юу вэ?}

Expo нь React Native дээр суурилсан платформ бөгөөд гар утасны аппликейшн хөгжүүлэлтийг илүү хялбар болгодог. Expo нь "managed workflow" болон "bare workflow" гэсэн хоёр төрлийн хөгжүүлэлтийн загварыг санал болгодог.

\textbf{Expo-ийн давуу талууд:}
\begin{itemize}
    \item \textbf{Хялбар эхлэл:} Xcode (iOS) эсвэл Android Studio (Android) суулгахгүйгээр хөгжүүлэлт эхлүүлэх боломжтой. \texttt{npx create-expo-app} командаар шинэ төсөл үүсгэхэд 2 минут хангалттай.
    \item \textbf{Expo Go апп:} Бодит төхөөрөмж дээр шууд туршилт хийх боломжтой. QR код уншуулахад л хангалттай.
    \item \textbf{Бэлэн SDK:} Камер, байршил (GPS), мэдэгдэл (Push Notifications), файл систем, биометрик баталгаажуулалт зэрэг 50 гаруй бэлэн модуль байдаг. Энэ нь native код бичих шаардлагагүй болгодог.
    \item \textbf{OTA Updates (Over-The-Air):} App Store/Google Play-д дахин илгээхгүйгээр JavaScript кодын шинэчлэлтийг хэрэглэгчдэд шууд хүргэх боломжтой. Энэ нь засвар хийх хугацааг хэдэн өдрөөс хэдэн минут болгон багасгадаг.
    \item \textbf{EAS Build:} Үүлэн дээр build хийх боломжтой тул локал машинд Xcode/Android Studio суулгах шаардлагагүй.
    \item \textbf{Expo Router:} Файлын бүтцэд суурилсан навигаци систем нь Next.js-тэй төстэй routing боломжийг санал болгодог.
\end{itemize}

\textbf{React Native Web}

React Native Web нь React Native компонентуудыг вэб хөтөч дээр ажиллуулах боломжийг олгодог. Энэ нь нэг кодоор гар утас (iOS, Android) болон вэб аппликейшн хөгжүүлэх боломжийг бүрдүүлдэг. Twitter, Uber Eats зэрэг томоохон компаниуд энэ аргачлалыг ашигладаг.

\subsubsection{Database ба Backend (Firebase)}
Firebase нь Google-ийн үүлэн тооцооллын платформ бөгөөд гар утас болон вэб аппликейшн хөгжүүлэхэд зориулагдсан цогц шийдэл юм. "Backend-as-a-Service" (BaaS) загвараар ажилладаг бөгөөд энэ нь сервер талын инфраструктурыг бүрэн удирдах шаардлагагүй болгодог.

\begin{table}[H]
\centering
\caption{Firebase үйлчилгээнүүд}
\label{tab:firebase-services}
\begin{tabular}{|l|p{9cm}|}
\hline
\textbf{Үйлчилгээ} & \textbf{Зориулалт} \\
\hline
Firebase (v12.4.0) & Backend-as-a-Service платформ \\
\hline
Cloud Firestore & NoSQL өгөгдлийн сан, бодит цагийн синхрончлол \\
\hline
Firebase Authentication & Хэрэглэгчийн баталгаажуулалт (Утас, Email) \\
\hline
Firebase Storage & Файл хадгалах (зураг гэх мэт) \\
\hline
Firebase Realtime Database & Бодит цагийн өгөгдөл \\
\hline
\end{tabular}
\end{table}

\textbf{Firebase гэж юу вэ?}

Firebase нь анх 2011 онд стартап компани болон байгуулагдсан бөгөөд 2014 онд Google компани худалдан авсан. Өнөөдөр Firebase нь Google Cloud Platform-ийн нэг хэсэг бөгөөд дэлхий даяар 3 сая гаруй апп ашигладаг. Firebase нь "Serverless" архитектурыг дэмждэг тул хөгжүүлэгч сервер удирдах, scale хийх зэрэг асуудлуудаар толгой өвдөхгүй.

\textbf{Cloud Firestore - NoSQL Өгөгдлийн сан}

Cloud Firestore нь Firebase-ийн хамгийн чухал үйлчилгээ бөгөөд NoSQL документ-суурьтай (document-oriented) өгөгдлийн сан юм. Firestore нь өгөгдлийг collection (цуглуулга) болон document (баримт) бүтцээр хадгалдаг.

\textbf{Firestore-ийн өгөгдлийн бүтэц:}
\begin{lstlisting}[language=JavaScript]
// Firestore өгөгдлийн бүтцийн жишээ
orders (collection)
  |-- orderId1 (document)
  |     |-- customerId: "cust123"
  |     |-- items: [{name: "Product A", qty: 10}]
  |     |-- status: "pending"
  |     |-- createdAt: Timestamp
  |-- orderId2 (document)
        |-- ...
\end{lstlisting}

\textbf{Cloud Firestore-ийн онцлог:}
\begin{itemize}
    \item \textbf{Real-time синхрончлол:} Өгөгдөл өөрчлөгдөхөд бүх холбогдсон клиентүүдэд шууд мэдэгддэг. \texttt{onSnapshot()} функцээр бодит цагийн сонсогч (listener) тохируулж болно. Энэ нь захиалгын төлөв өөрчлөгдөхөд хэрэглэгчид шууд мэдэгдэх боломжийг олгодог.
    \item \textbf{Offline Persistence (Оффлайн дэмжлэг):} Интернэт холболтгүй үед ч аппликейшн ажиллах чадвартай. Өгөгдөл локал кэшэд хадгалагдаж, холболт сэргэхэд автоматаар синхрончлогддог. Энэ нь хүргэлтийн жолооч интернэтгүй газарт явах үед маш чухал.
    \item \textbf{Scalability (Өргөтгөх чадвар):} Хэрэглэгчийн тоо нэмэгдэхэд автоматаар scale хийгддэг. Google-ийн дэд бүтцийг ашигладаг тул сая сая хэрэглэгчид үйлчлэх боломжтой.
    \item \textbf{Security Rules (Аюулгүй байдлын дүрмүүд):} Нарийвчилсан хандалтын эрхийн удирдлага. Хэн ямар өгөгдөлд хандаж болохыг тодорхойлох боломжтой. Жишээ нь: хэрэглэгч зөвхөн өөрийн захиалгыг үзэх боломжтой.
    \item \textbf{Query боломжууд:} Compound query (олон нөхцөлтэй хайлт), ordering, pagination зэрэг боломжууд бий. Гэхдээ SQL шиг JOIN үйлдэл хийх боломжгүй.
\end{itemize}

\textbf{Firebase Authentication - Баталгаажуулалт}

Firebase Authentication нь хэрэглэгчийн бүртгэл, нэвтрэлтийг удирдах бэлэн шийдэл юм. Дараах аргуудыг дэмждэг:
\begin{itemize}
    \item \textbf{Имэйл/Нууц үг:} Уламжлалт бүртгэлийн арга
    \item \textbf{Утасны дугаар (SMS):} OTP код илгээж баталгаажуулах. Монголд маш түгээмэл.
    \item \textbf{Social Login:} Google, Facebook, Apple, Twitter зэрэг нийгмийн сүлжээгээр нэвтрэх
    \item \textbf{Anonymous:} Бүртгэлгүйгээр түр зуурын хэрэглэгч үүсгэх
\end{itemize}

Tdelivery системд утасны дугаараар баталгаажуулалт (Phone Authentication) ашигласан. Энэ нь Монголын хэрэглэгчдэд илүү тохиромжтой - имэйл хаяг байнга ашигладаггүй хүмүүст ч хялбар.

\textbf{Firebase Storage - Файл хадгалалт}

Firebase Storage нь зураг, видео, аудио зэрэг файлуудыг үүлэн дээр хадгалах үйлчилгээ юм. Google Cloud Storage дээр суурилсан тул өндөр найдвартай, хурдан байдаг. Tdelivery системд дараах зориулалтаар ашигласан:
\begin{itemize}
    \item Бүтээгдэхүүний зураг хадгалах
    \item Хүргэлтийн баталгаажуулалтын зураг (e-POD)
    \item Хэрэглэгчийн профайл зураг
\end{itemize}

\textbf{Firebase-ийн үнийн бодлого}

Firebase нь "Spark" (үнэгүй) болон "Blaze" (pay-as-you-go) гэсэн хоёр үнийн төлөвлөгөөтэй. Spark төлөвлөгөө нь жижиг төсөл, сурах зориулалтаар хангалттай:
\begin{itemize}
    \item Firestore: 1GB хадгалалт, 50,000 унших, 20,000 бичих өдөрт
    \item Authentication: Хязгааргүй хэрэглэгч
    \item Storage: 5GB хадгалалт
\end{itemize}

\subsubsection{NoSQL ба Atomic Denormalization}
Уламжлалт RDBMS (Relational Database Management System) -ээс ялгаатай нь NoSQL өгөгдлийн сан нь өгөгдлийг баримт (Document) хэлбэрээр хадгалдаг. Энэ хэсэгт NoSQL-ийн онцлог, давуу тал, сул тал болон Firestore-д ашигладаг denormalization аргачлалыг дэлгэрэнгүй авч үзнэ.

\textbf{SQL vs NoSQL харьцуулалт}

\begin{table}[H]
\centering
\caption{SQL ба NoSQL харьцуулалт}
\label{tab:sql-nosql}
\begin{tabular}{|l|p{5.5cm}|p{5.5cm}|}
\hline
\textbf{Шалгуур} & \textbf{SQL (RDBMS)} & \textbf{NoSQL (Document DB)} \\
\hline
Өгөгдлийн бүтэц & Хүснэгт, мөр, багана (Table, Row, Column) & Collection, Document, Field \\
\hline
Схем & Хатуу схемтэй (Rigid Schema) & Уян хатан схем (Flexible Schema) \\
\hline
Холбоос & JOIN үйлдлээр холбогдсон хүснэгтүүд & Embedded documents, References \\
\hline
Scale & Вертикал (илүү хүчтэй сервер) & Хоризонтал (илүү олон сервер) \\
\hline
ACID & Бүрэн дэмждэг & Хязгаарлагдмал (document түвшинд) \\
\hline
Хэрэглээ & Банк, санхүү, нарийн холбоостой өгөгдөл & Вэб апп, гар утас, бодит цагийн апп \\
\hline
\end{tabular}
\end{table}

\textbf{NoSQL-ийн давуу талууд гар утасны аппд:}
\begin{itemize}
    \item \textbf{Уян хатан бүтэц:} Схем өөрчлөхөд migration хийх шаардлагагүй. Шинэ field нэмэхэд хуучин өгөгдөлд нөлөөлөхгүй.
    \item \textbf{Бодит цагийн синхрончлол:} WebSocket-ээр өгөгдөл шууд push хийгддэг.
    \item \textbf{Оффлайн дэмжлэг:} Локал кэш, conflict resolution бэлэн.
    \item \textbf{JSON формат:} JavaScript/TypeScript-тэй шууд нийцтэй.
\end{itemize}

\textbf{Atomic Denormalization гэж юу вэ?}

Firebase Firestore-д өгөгдлийн унших хурдыг нэмэгдүүлэхийн тулд "Denormalization" аргыг ашигладаг. Энэ нь нэг өгөгдлийг олон газар хувилж хадгалах арга бөгөөд JOIN үйлдэл хийх шаардлагагүй болгож, query-ийн хурдыг эрс нэмэгдүүлдэг.

\textbf{Denormalization-ийн жишээ:}

SQL-д захиалгын мэдээлэл авахад 3 хүснэгтийг JOIN хийх шаардлагатай:
\begin{lstlisting}[language=SQL]
-- SQL: 3 хүснэгт JOIN хийх
SELECT o.*, c.name, c.phone, p.name as product_name
FROM orders o
JOIN customers c ON o.customer_id = c.id
JOIN products p ON o.product_id = p.id
WHERE o.id = 'order123';
\end{lstlisting}

NoSQL (Firestore)-д denormalization ашиглан нэг document-д бүх мэдээллийг хадгална:
\begin{lstlisting}[language=JavaScript]
// Firestore: Denormalized document
{
  "orderId": "order123",
  "customer": {
    "id": "cust456",
    "name": "Болд",        // Хуулбарласан
    "phone": "99001122"    // Хуулбарласан
  },
  "product": {
    "id": "prod789",
    "name": "Coca-Cola 1.5L", // Хуулбарласан
    "price": 3500
  },
  "quantity": 10,
  "totalAmount": 35000
}
\end{lstlisting}

\textbf{Atomic Update - Өгөгдлийн нэгдмэл байдал хангах}

Denormalization-ийн сул тал нь өгөгдөл олон газар хуулбарлагдсан тул нэг газар өөрчлөхөд бүх газар өөрчлөх шаардлагатай. Үүнийг "Atomic Update" буюу batch/transaction ашиглан шийддэг:

\begin{lstlisting}[language=JavaScript]
// Batch write - Бүх өөрчлөлтийг нэг дор хийх
const batch = writeBatch(db);

// Үндсэн бүтээгдэхүүний нэр өөрчлөх
batch.update(doc(db, 'products', 'prod789'), {
  name: 'Coca-Cola 2L'
});

// Бүх холбогдох захиалгуудын нэрийг өөрчлөх
orderIds.forEach(orderId => {
  batch.update(doc(db, 'orders', orderId), {
    'product.name': 'Coca-Cola 2L'
  });
});

await batch.commit(); // Бүгд амжилттай эсвэл бүгд буцаагдана
\end{lstlisting}

\textbf{Denormalization-ийн зарчмууд:}
\begin{enumerate}
    \item \textbf{Унших давтамж > Бичих давтамж:} Ихэвчлэн уншигддаг өгөгдлийг denormalize хийх нь илүү үр ашигтай.
    \item \textbf{Хамтдаа хэрэглэгддэг өгөгдөл:} Нэг дэлгэц дээр харуулах өгөгдлийг нэг document-д хадгалах.
    \item \textbf{Бага өөрчлөгддөг өгөгдөл:} Хэрэглэгчийн нэр, утас зэрэг бага өөрчлөгддөг өгөгдлийг хуулбарлах нь аюулгүй.
\end{enumerate}

\subsubsection{Навигаци (Navigation)}
Аппликейшний дэлгэцүүдийн хооронд шилжихэд React Navigation санг ашигласан. Навигаци нь гар утасны аппликейшний хэрэглэгчийн туршлагын (UX) чухал хэсэг бөгөөд зөв зохион байгуулалт нь аппын ашиглахад хялбар байдлыг тодорхойлдог.

\begin{table}[H]
\centering
\caption{Навигацийн сангууд}
\label{tab:navigation}
\begin{tabular}{|l|l|p{6cm}|}
\hline
\textbf{Сан} & \textbf{Хувилбар} & \textbf{Зориулалт} \\
\hline
@react-navigation/native & \textasciicircum7.1.8 & Үндсэн навигаци \\
\hline
@react-navigation/native-stack & \textasciicircum7.1.9 & Stack навигаци (дэлгэц дээр дэлгэц нэмэх) \\
\hline
@react-navigation/bottom-tabs & \textasciicircum7.4.0 & Доод таб навигаци \\
\hline
\end{tabular}
\end{table}

\textbf{React Navigation гэж юу вэ?}

React Navigation нь React Native аппликейшнд навигаци хийхэд хамгийн өргөн хэрэглэгддэг сан юм. Энэ нь JavaScript дээр бичигдсэн тул cross-platform бөгөөд iOS, Android хоёуланд нэг кодоор ажилладаг.

\textbf{Навигацийн төрлүүд:}

\begin{itemize}
    \item \textbf{Stack Navigator:} Дэлгэцүүдийг "stack" (овоолго) шиг удирддаг. Шинэ дэлгэц нээхэд дээр нь нэмэгдэж, буцах товч дарахад дээрээс хасагдана. Ихэнх апп-ийн үндсэн навигаци.
    
    \item \textbf{Bottom Tab Navigator:} Дэлгэцийн доод хэсэгт байрлах табууд. Үндсэн функцуудад хурдан хандах боломжтой. Tdelivery системд: Нүүр, Захиалга, Сагс, Профайл гэсэн 4 таб ашигласан.
    
    \item \textbf{Drawer Navigator:} Хажуугаас гулсаж гарах цэс. Олон үйлдэлтэй аппуудад тохиромжтой.
    
    \item \textbf{Top Tab Navigator:} Дээд хэсэгт байрлах табууд. Material Design-д түгээмэл.
\end{itemize}

\textbf{Tdelivery системийн навигацийн бүтэц:}

\begin{lstlisting}[language=JavaScript]
// Навигацийн бүтэц
RootNavigator
  |-- AuthStack (Нэвтрээгүй үед)
  |     |-- LoginScreen
  |     |-- RegisterScreen
  |     |-- OTPVerificationScreen
  |
  |-- MainTabs (Нэвтэрсэн үед)
        |-- HomeStack
        |     |-- HomeScreen
        |     |-- ProductDetailScreen
        |     |-- CategoryScreen
        |
        |-- OrdersStack
        |     |-- OrderListScreen
        |     |-- OrderDetailScreen
        |
        |-- CartStack
        |     |-- CartScreen
        |     |-- CheckoutScreen
        |
        |-- ProfileStack
              |-- ProfileScreen
              |-- EditProfileScreen
              |-- SettingsScreen
\end{lstlisting}

\textbf{Navigation-ийн чухал ойлголтууд:}
\begin{itemize}
    \item \textbf{Screen:} Бие даасан дэлгэц, компонент
    \item \textbf{Route:} Дэлгэц рүү очих зам, параметрүүд
    \item \textbf{Navigator:} Дэлгэцүүдийг удирдах контейнер
    \item \textbf{Navigation Prop:} Дэлгэц хооронд шилжих функцууд (\texttt{navigate}, \texttt{goBack}, \texttt{push})
    \item \textbf{Route Prop:} Дэлгэцэнд дамжуулсан параметрүүд (\texttt{route.params})
\end{itemize}

\subsubsection{UI/UX Сангууд}
Хэрэглэгчийн интерфэйсийг сайжруулахад дараах сангуудыг ашигласан. UI (User Interface) нь хэрэглэгчийн харж буй визуал элементүүд, UX (User Experience) нь хэрэглэгчийн аппыг ашиглах туршлага юм.

\begin{table}[H]
\centering
\caption{UI/UX сангууд}
\label{tab:ui-libraries}
\begin{tabular}{|l|l|p{6cm}|}
\hline
\textbf{Сан} & \textbf{Хувилбар} & \textbf{Зориулалт} \\
\hline
React Native Paper & \textasciicircum5.12.3 & Material Design UI компонентууд \\
\hline
@expo/vector-icons & \textasciicircum15.0.3 & Дүрс тэмдэг (Icons) \\
\hline
NativeWind & \textasciicircum4.2.1 & Tailwind CSS React Native-д \\
\hline
TailwindCSS & \textasciicircum3.4.18 & Utility-first CSS фреймворк \\
\hline
Expo Image & \textasciitilde3.0.10 & Зураг харуулах, кэшлэх \\
\hline
\end{tabular}
\end{table}

\textbf{React Native Paper}

React Native Paper нь Google-ийн Material Design дизайны хэлийг React Native-д хэрэгжүүлсэн UI компонент сан юм. 30 гаруй бэлэн компонент агуулдаг:

\begin{itemize}
    \item \textbf{Button:} Primary, outlined, text, contained төрлүүдтэй товч
    \item \textbf{Card:} Бүтээгдэхүүн, захиалга харуулахад тохиромжтой карт
    \item \textbf{TextInput:} Текст оруулах талбар, validation дэмжлэгтэй
    \item \textbf{Snackbar:} Түр зуурын мэдэгдэл (toast)
    \item \textbf{Dialog:} Модал цонх
    \item \textbf{FAB (Floating Action Button):} Хөвөгч үйлдлийн товч
    \item \textbf{ActivityIndicator:} Ачаалж байгааг харуулах индикатор
\end{itemize}

\textbf{Material Design-ийн давуу талууд:}
\begin{itemize}
    \item Google-ийн дизайны стандарт тул хэрэглэгчдэд танил
    \item Accessibility (хүртээмжтэй байдал) бэлэн хийгдсэн
    \item Theming (өнгө, фонт тохируулах) боломжтой
    \item iOS болон Android дээр нэгдмэл харагдах
\end{itemize}

\textbf{NativeWind ба TailwindCSS}

TailwindCSS нь "utility-first" CSS фреймворк бөгөөд NativeWind нь үүнийг React Native-д ашиглах боломжийг олгодог. Уламжлалт CSS-ээс ялгаатай нь бэлэн class-уудыг шууд ашиглана.

\textbf{Уламжлалт StyleSheet vs NativeWind харьцуулалт:}

\begin{lstlisting}[language=JavaScript]
// Уламжлалт React Native StyleSheet
const styles = StyleSheet.create({
  container: {
    padding: 16,
    backgroundColor: 'white',
    borderRadius: 8,
    shadowColor: '#000',
    shadowOffset: { width: 0, height: 2 },
    shadowOpacity: 0.1,
  },
  title: {
    fontSize: 18,
    fontWeight: 'bold',
    color: '#1f2937',
  }
});

<View style={styles.container}>
  <Text style={styles.title}>Захиалга</Text>
</View>
\end{lstlisting}

\begin{lstlisting}[language=JavaScript]
// NativeWind (TailwindCSS)
<View className="p-4 bg-white rounded-lg shadow-md">
  <Text className="text-lg font-bold text-gray-800">Захиалга</Text>
</View>
\end{lstlisting}

\textbf{NativeWind/TailwindCSS-ийн давуу талууд:}
\begin{itemize}
    \item \textbf{Utility-first:} Бэлэн class-уудаар хурдан стиль хийх боломжтой. \texttt{p-4} нь padding: 16px, \texttt{bg-white} нь цагаан дэвсгэр.
    \item \textbf{Responsive:} Дэлгэцийн хэмжээнд тохируулсан стиль. \texttt{md:p-8} нь дунд дэлгэц дээр илүү том padding.
    \item \textbf{Consistency:} Урьдчилан тодорхойлсон утгууд (spacing, colors) нь дизайны нэгдмэл байдлыг хангах.
    \item \textbf{Жижиг bundle:} Зөвхөн ашигласан стилүүд bundle-д орно (PurgeCSS).
    \item \textbf{Dark mode:} \texttt{dark:bg-gray-900} гэх мэтээр хар горимыг хялбар тохируулах.
\end{itemize}

\textbf{@expo/vector-icons}

Expo Vector Icons нь 10+ icon set-ийг агуулсан сан юм:
\begin{itemize}
    \item \textbf{MaterialIcons:} Google-ийн Material Design icons
    \item \textbf{FontAwesome:} Хамгийн түгээмэл icon сан
    \item \textbf{Ionicons:} iOS-д зориулсан icons
    \item \textbf{Feather:} Хялбар, цэвэрхэн icons
\end{itemize}

\begin{lstlisting}[language=JavaScript]
import { MaterialIcons, FontAwesome } from '@expo/vector-icons';

<MaterialIcons name="shopping-cart" size={24} color="blue" />
<FontAwesome name="user" size={20} color="#333" />
\end{lstlisting}

\textbf{Expo Image}

Expo Image нь React Native-ийн үндсэн Image компонентоос илүү хүчирхэг:
\begin{itemize}
    \item \textbf{Automatic caching:} Зураг автоматаар кэшлэгдэнэ
    \item \textbf{Progressive loading:} Бүдэг зургаас тод зураг руу шилжих
    \item \textbf{Placeholder:} Ачаалж байх үеийн түр зураг
    \item \textbf{BlurHash:} Зургийн өнгөний товч дүрслэл
    \item \textbf{WebP support:} Шахагдсан зургийн формат дэмжлэг
\end{itemize}

\subsubsection{Expo SDK Сангууд}
Expo платформын бэлэн сангуудыг ашиглан төхөөрөмжийн боломжуудад хандсан. Expo SDK нь native код бичихгүйгээр төхөөрөмжийн hardware болон software боломжуудад хандах боломжийг олгодог.

\begin{table}[H]
\centering
\caption{Expo SDK сангууд}
\label{tab:expo-sdk}
\begin{tabular}{|l|l|p{6cm}|}
\hline
\textbf{Сан} & \textbf{Хувилбар} & \textbf{Зориулалт} \\
\hline
expo-constants & \textasciitilde18.0.10 & Апп-ийн тогтмол утгууд \\
\hline
expo-font & \textasciitilde14.0.9 & Фонт удирдлага \\
\hline
expo-haptics & \textasciitilde15.0.7 & Чичиргээ (Haptic feedback) \\
\hline
expo-image-picker & \textasciicircum17.0.8 & Зураг сонгох \\
\hline
expo-linking & \textasciitilde8.0.8 & Deep linking \\
\hline
expo-splash-screen & \textasciitilde31.0.10 & Splash дэлгэц \\
\hline
expo-status-bar & \textasciitilde3.0.8 & Статус бар удирдлага \\
\hline
expo-web-browser & \textasciicircum15.0.8 & Веб хөтөч нээх \\
\hline
\end{tabular}
\end{table}

\textbf{Expo SDK сангуудын дэлгэрэнгүй тайлбар:}

\textbf{expo-constants} - Апп болон төхөөрөмжийн мэдээлэлд хандах:
\begin{lstlisting}[language=JavaScript]
import Constants from 'expo-constants';

// Апп-ийн хувилбар
console.log(Constants.expoConfig.version); // "1.0.0"
// Төхөөрөмжийн нэр
console.log(Constants.deviceName); // "iPhone 14"
// Платформ
console.log(Constants.platform); // { ios: {...} }
\end{lstlisting}

\textbf{expo-font} - Custom фонт ачаалах:
\begin{lstlisting}[language=JavaScript]
import { useFonts } from 'expo-font';

const [fontsLoaded] = useFonts({
  'Roboto-Bold': require('./assets/fonts/Roboto-Bold.ttf'),
  'Roboto-Regular': require('./assets/fonts/Roboto-Regular.ttf'),
});

// Фонт ашиглах
<Text style={{ fontFamily: 'Roboto-Bold' }}>Текст</Text>
\end{lstlisting}

\textbf{expo-haptics} - Чичиргээ (Haptic feedback):

Хэрэглэгч товч дарах, swipe хийх зэрэг үйлдэл хийхэд физик мэдрэмж өгөх:
\begin{lstlisting}[language=JavaScript]
import * as Haptics from 'expo-haptics';

// Товч дарахад хөнгөн чичиргээ
const handlePress = () => {
  Haptics.impactAsync(Haptics.ImpactFeedbackStyle.Light);
};

// Алдаа гарахад хүчтэй чичиргээ
const handleError = () => {
  Haptics.notificationAsync(Haptics.NotificationFeedbackType.Error);
};

// Захиалга амжилттай болоход
const handleSuccess = () => {
  Haptics.notificationAsync(Haptics.NotificationFeedbackType.Success);
};
\end{lstlisting}

\textbf{expo-image-picker} - Зураг сонгох:

Галерей эсвэл камераас зураг авах. Tdelivery системд бүтээгдэхүүний зураг, хүргэлтийн баталгаажуулалтын зураг авахад ашигласан:
\begin{lstlisting}[language=JavaScript]
import * as ImagePicker from 'expo-image-picker';

const pickImage = async () => {
  // Зөвшөөрөл авах
  const { status } = await ImagePicker.requestMediaLibraryPermissionsAsync();
  
  if (status !== 'granted') {
    alert('Галерейд хандах зөвшөөрөл шаардлагатай');
    return;
  }
  
  // Зураг сонгох
  const result = await ImagePicker.launchImageLibraryAsync({
    mediaTypes: ImagePicker.MediaTypeOptions.Images,
    allowsEditing: true,
    aspect: [4, 3],
    quality: 0.8,
  });
  
  if (!result.canceled) {
    console.log(result.assets[0].uri);
  }
};
\end{lstlisting}

\textbf{expo-linking} - Deep linking:

Бусад апп эсвэл вэб хуудаснаас аппыг нээх, тодорхой дэлгэц рүү шилжих:
\begin{lstlisting}[language=JavaScript]
import * as Linking from 'expo-linking';

// Апп дотроос утас руу залгах
Linking.openURL('tel:+97699001122');

// Email илгээх
Linking.openURL('mailto:support@tdelivery.mn');

// Google Maps нээх
Linking.openURL('https://maps.google.com/?q=47.9184,106.9177');

// Deep link үүсгэх
const url = Linking.createURL('order/12345');
// tdelivery://order/12345
\end{lstlisting}

\textbf{expo-splash-screen} - Splash дэлгэц:

Апп ачаалж байх үед харуулах эхний дэлгэц:
\begin{lstlisting}[language=JavaScript]
import * as SplashScreen from 'expo-splash-screen';

// Splash дэлгэцийг байлгах
SplashScreen.preventAutoHideAsync();

// Өгөгдөл ачаалсны дараа нуух
useEffect(() => {
  async function prepare() {
    await loadFonts();
    await fetchInitialData();
    await SplashScreen.hideAsync();
  }
  prepare();
}, []);
\end{lstlisting}

\textbf{expo-status-bar} - Статус бар удирдлага:

Дэлгэцийн дээд хэсэгт байрлах статус бар (цаг, батарей, сүлжээ)-ын харагдах байдлыг удирдах:
\begin{lstlisting}[language=JavaScript]
import { StatusBar } from 'expo-status-bar';

// Хар дэвсгэр дээр цагаан текст
<StatusBar style="light" />

// Цагаан дэвсгэр дээр хар текст
<StatusBar style="dark" />

// Автомат (дэвсгэрээс хамаарч)
<StatusBar style="auto" />
\end{lstlisting}

\subsubsection{State Management ба Storage}
Аппликейшний төлөв удирдлага болон локал хадгалалтад дараах шийдлүүдийг ашигласан. State Management нь аппликейшний өгөгдлийг (төлвийг) компонентууд хооронд хуваалцах, синхрончлох үйл явц юм.

\begin{table}[H]
\centering
\caption{State Management ба Storage}
\label{tab:state-management}
\begin{tabular}{|l|l|p{6cm}|}
\hline
\textbf{Сан} & \textbf{Хувилбар} & \textbf{Зориулалт} \\
\hline
React Context API & - & Төлөв удирдлага (AuthContext, CartContext) \\
\hline
@react-native-async-storage/async-storage & \textasciicircum2.2.0 & Локал хадгалалт \\
\hline
\end{tabular}
\end{table}

\textbf{State Management гэж юу вэ?}

Гар утасны аппликейшнд олон төрлийн өгөгдөл (төлөв) байдаг:
\begin{itemize}
    \item \textbf{Local State:} Нэг компонентын дотоод төлөв (жишээ нь: form input-ийн утга)
    \item \textbf{Global State:} Олон компонентод хэрэглэгддэг төлөв (жишээ нь: нэвтэрсэн хэрэглэгчийн мэдээлэл)
    \item \textbf{Server State:} Серверээс ирсэн өгөгдөл (жишээ нь: бүтээгдэхүүний жагсаалт)
    \item \textbf{URL State:} Navigation параметрүүд
\end{itemize}

\textbf{React Context API}

React Context API нь React-ийн суурь боломж бөгөөд global state удирдахад хялбар шийдэл юм. Redux, MobX зэрэг гуравдагч талын сангуудаас хялбар, жижиг-дунд хэмжээний аппуудад тохиромжтой.

\textbf{Context-ийн ажиллах зарчим:}
\begin{enumerate}
    \item \textbf{Context үүсгэх:} \texttt{createContext()} функцээр context үүсгэнэ
    \item \textbf{Provider:} Context-ийн утгыг хангах компонент
    \item \textbf{Consumer:} Context-ийн утгыг ашиглах компонент (\texttt{useContext} hook)
\end{enumerate}

\textbf{Tdelivery системийн Context-үүд:}

\begin{lstlisting}[language=JavaScript]
// AuthContext - Хэрэглэгчийн баталгаажуулалт
const AuthContext = createContext<AuthContextType | null>(null);

interface AuthContextType {
  user: User | null;
  isLoading: boolean;
  signIn: (phone: string, otp: string) => Promise<void>;
  signOut: () => Promise<void>;
}

export const AuthProvider: React.FC = ({ children }) => {
  const [user, setUser] = useState<User | null>(null);
  const [isLoading, setIsLoading] = useState(true);
  
  // Firebase Auth-тай холбогдох
  useEffect(() => {
    const unsubscribe = onAuthStateChanged(auth, (firebaseUser) => {
      setUser(firebaseUser);
      setIsLoading(false);
    });
    return unsubscribe;
  }, []);
  
  const signIn = async (phone: string, otp: string) => {
    // Нэвтрэх логик
  };
  
  const signOut = async () => {
    await auth.signOut();
  };
  
  return (
    <AuthContext.Provider value={{ user, isLoading, signIn, signOut }}>
      {children}
    </AuthContext.Provider>
  );
};

// Hook ашиглах
export const useAuth = () => useContext(AuthContext);
\end{lstlisting}

\begin{lstlisting}[language=JavaScript]
// CartContext - Сагсны удирдлага
const CartContext = createContext<CartContextType | null>(null);

interface CartItem {
  productId: string;
  name: string;
  price: number;
  quantity: number;
}

interface CartContextType {
  items: CartItem[];
  addItem: (product: Product) => void;
  removeItem: (productId: string) => void;
  updateQuantity: (productId: string, quantity: number) => void;
  clearCart: () => void;
  totalAmount: number;
  itemCount: number;
}
\end{lstlisting}

\textbf{React Context API-ийн давуу талууд:}
\begin{itemize}
    \item \textbf{Суурь сан:} Нэмэлт сан суулгах шаардлагагүй, React-д бэлэн байдаг.
    \item \textbf{Prop drilling шийдэл:} Олон түвшний компонентуудаар дамжуулахгүйгээр өгөгдөл дамжуулах. Жишээ нь: App -> Header -> UserMenu гэсэн бүтцэд Header-ээр user дамжуулах шаардлагагүй.
    \item \textbf{Хялбар:} Redux-ээс хялбар, boilerplate код бага. Redux-д 5-6 файл үүсгэх шаардлагатай бол Context-д 1 файл хангалттай.
    \item \textbf{TypeScript дэмжлэг:} Төрлийн шалгалт сайн ажилладаг.
\end{itemize}

\textbf{AsyncStorage - Локал хадгалалт}

AsyncStorage нь key-value хосоор өгөгдөл хадгалах локал хадгалалт юм. Вэб хөтчийн localStorage-тай төстэй боловч асинхрон ажилладаг.

\textbf{AsyncStorage-ийн хэрэглээ:}
\begin{lstlisting}[language=JavaScript]
import AsyncStorage from '@react-native-async-storage/async-storage';

// Өгөгдөл хадгалах
await AsyncStorage.setItem('userToken', 'abc123');
await AsyncStorage.setItem('settings', JSON.stringify({
  notifications: true,
  darkMode: false
}));

// Өгөгдөл унших
const token = await AsyncStorage.getItem('userToken');
const settings = JSON.parse(await AsyncStorage.getItem('settings'));

// Өгөгдөл устгах
await AsyncStorage.removeItem('userToken');

// Бүх өгөгдөл устгах
await AsyncStorage.clear();
\end{lstlisting}

\textbf{AsyncStorage-ийн хэрэглээний жишээ Tdelivery системд:}
\begin{itemize}
    \item Хэрэглэгчийн тохиргоо (мэдэгдэл, хар горим)
    \item Сагсны түр өгөгдөл
    \item Сүүлд үзсэн бүтээгдэхүүнүүд
    \item Онбординг харуулсан эсэх
\end{itemize}

\subsubsection{Network ба API}
Сервертэй харилцахад дараах сангуудыг ашигласан. Network layer нь аппликейшн болон серверийн хоорондын харилцааг удирддаг чухал хэсэг юм.

\begin{table}[H]
\centering
\caption{Network ба API}
\label{tab:network-api}
\begin{tabular}{|l|l|p{6cm}|}
\hline
\textbf{Сан} & \textbf{Хувилбар} & \textbf{Зориулалт} \\
\hline
Axios & \textasciicircum1.7.0 & HTTP хүсэлт илгээх \\
\hline
Firebase SDK & \textasciicircum12.4.0 & Firebase API-тай харилцах \\
\hline
\end{tabular}
\end{table}

\textbf{Axios гэж юу вэ?}

Axios нь JavaScript-д зориулагдсан HTTP client сан бөгөөд REST API-тай харилцахад хамгийн түгээмэл ашиглагддаг. Fetch API-аас илүү олон боломжтой.

\textbf{Axios-ийн давуу талууд:}
\begin{itemize}
    \item \textbf{Автомат JSON:} Response-ийг автоматаар JSON болгон parse хийнэ
    \item \textbf{Interceptors:} Request/Response-ийг дунд нь барьж өөрчлөх боломжтой
    \item \textbf{Timeout:} Хүсэлтийн хугацаа хязгаарлах
    \item \textbf{Cancel tokens:} Хүсэлтийг цуцлах боломжтой
    \item \textbf{Error handling:} Алдааны мэдээлэл дэлгэрэнгүй
\end{itemize}

\textbf{Axios-ийн тохиргоо:}
\begin{lstlisting}[language=JavaScript]
import axios from 'axios';

// Axios instance үүсгэх
const api = axios.create({
  baseURL: 'https://api.tdelivery.mn/v1',
  timeout: 10000, // 10 секунд
  headers: {
    'Content-Type': 'application/json',
  },
});

// Request interceptor - токен нэмэх
api.interceptors.request.use(
  async (config) => {
    const token = await AsyncStorage.getItem('authToken');
    if (token) {
      config.headers.Authorization = `Bearer ${token}`;
    }
    return config;
  },
  (error) => Promise.reject(error)
);

// Response interceptor - алдаа боловсруулах
api.interceptors.response.use(
  (response) => response.data,
  (error) => {
    if (error.response?.status === 401) {
      // Токен хүчингүй - дахин нэвтрүүлэх
      navigation.navigate('Login');
    }
    return Promise.reject(error);
  }
);
\end{lstlisting}

\textbf{Firebase SDK}

Firebase SDK нь Firebase үйлчилгээнүүдтэй (Firestore, Auth, Storage) холбогдох албан ёсны сан юм. REST API биш, WebSocket-ээр бодит цагийн холболт үүсгэдэг.

\textbf{Firebase SDK-ийн хэрэглээ:}
\begin{lstlisting}[language=JavaScript]
import { initializeApp } from 'firebase/app';
import { getFirestore, collection, getDocs, onSnapshot } from 'firebase/firestore';
import { getAuth } from 'firebase/auth';

// Firebase тохиргоо
const firebaseConfig = {
  apiKey: "AIza...",
  authDomain: "tdelivery.firebaseapp.com",
  projectId: "tdelivery",
  storageBucket: "tdelivery.appspot.com",
  messagingSenderId: "123456789",
  appId: "1:123456789:web:abc123"
};

const app = initializeApp(firebaseConfig);
const db = getFirestore(app);
const auth = getAuth(app);

// Бүтээгдэхүүн авах (нэг удаа)
const getProducts = async () => {
  const snapshot = await getDocs(collection(db, 'products'));
  return snapshot.docs.map(doc => ({ id: doc.id, ...doc.data() }));
};

// Бодит цагийн сонсогч
const subscribeToOrders = (userId: string, callback: Function) => {
  return onSnapshot(
    query(collection(db, 'orders'), where('userId', '==', userId)),
    (snapshot) => {
      const orders = snapshot.docs.map(doc => ({ id: doc.id, ...doc.data() }));
      callback(orders);
    }
  );
};
\end{lstlisting}

\subsubsection{Animation ба Gestures}
Анимаци болон хэрэглэгчийн жестүүдийг удирдахад дараах сангуудыг ашигласан. Сайн анимаци нь хэрэглэгчийн туршлагыг эрс сайжруулж, аппыг илүү "амьд" мэт мэдрүүлдэг.

\begin{table}[H]
\centering
\caption{Animation ба Gestures}
\label{tab:animation}
\begin{tabular}{|l|l|p{6cm}|}
\hline
\textbf{Сан} & \textbf{Хувилбар} & \textbf{Зориулалт} \\
\hline
react-native-reanimated & \textasciitilde4.1.1 & Өндөр гүйцэтгэлтэй анимаци \\
\hline
react-native-gesture-handler & \textasciitilde2.28.0 & Жест (Gesture) удирдлага \\
\hline
react-native-screens & \textasciitilde4.16.0 & Дэлгэцийн оновчлол \\
\hline
react-native-safe-area-context & \textasciitilde5.6.0 & Safe area удирдлага \\
\hline
\end{tabular}
\end{table}

\textbf{React Native Reanimated}

Reanimated нь React Native-д өндөр гүйцэтгэлтэй анимаци хийх сан юм. React Native-ийн үндсэн Animated API-аас ялгаатай нь UI thread дээр шууд ажилладаг тул 60 FPS анимаци хийх боломжтой.

\textbf{Reanimated-ийн онцлог:}
\begin{itemize}
    \item \textbf{Worklets:} JavaScript код UI thread дээр ажиллах
    \item \textbf{Shared Values:} JS болон UI thread хооронд хуваалцах утгууд
    \item \textbf{Animated Styles:} Анимацитай стилүүд
    \item \textbf{Layout Animations:} Компонент нэмэх/хасахад автомат анимаци
\end{itemize}

\textbf{Reanimated-ийн жишээ:}
\begin{lstlisting}[language=JavaScript]
import Animated, {
  useSharedValue,
  useAnimatedStyle,
  withSpring,
  withTiming,
} from 'react-native-reanimated';

const AnimatedCard = () => {
  const scale = useSharedValue(1);
  const opacity = useSharedValue(1);
  
  const animatedStyle = useAnimatedStyle(() => ({
    transform: [{ scale: scale.value }],
    opacity: opacity.value,
  }));
  
  const handlePress = () => {
    // Spring анимаци - bounce effect
    scale.value = withSpring(0.95, {}, () => {
      scale.value = withSpring(1);
    });
  };
  
  const fadeOut = () => {
    // Timing анимаци - гөлгөр
    opacity.value = withTiming(0, { duration: 300 });
  };
  
  return (
    <Animated.View style={[styles.card, animatedStyle]}>
      <Pressable onPress={handlePress}>
        <Text>Анимацитай карт</Text>
      </Pressable>
    </Animated.View>
  );
};
\end{lstlisting}

\textbf{React Native Gesture Handler}

Gesture Handler нь touch event-үүдийг native түвшинд боловсруулдаг тул илүү responsive, хурдан ажилладаг.

\textbf{Дэмждэг жестүүд:}
\begin{itemize}
    \item \textbf{Tap:} Товших (single, double, long press)
    \item \textbf{Pan:} Чирэх
    \item \textbf{Pinch:} Хоёр хуруугаар томруулах/жижигрүүлэх
    \item \textbf{Rotation:} Эргүүлэх
    \item \textbf{Fling:} Түрхэх (swipe)
\end{itemize}

\textbf{Swipe to Delete жишээ:}
\begin{lstlisting}[language=JavaScript]
import { Gesture, GestureDetector } from 'react-native-gesture-handler';

const SwipeableItem = ({ onDelete }) => {
  const translateX = useSharedValue(0);
  
  const panGesture = Gesture.Pan()
    .onUpdate((event) => {
      // Зүүн тийш чирэх
      translateX.value = Math.min(0, event.translationX);
    })
    .onEnd(() => {
      if (translateX.value < -100) {
        // Устгах
        runOnJS(onDelete)();
      } else {
        // Буцаах
        translateX.value = withSpring(0);
      }
    });
  
  const animatedStyle = useAnimatedStyle(() => ({
    transform: [{ translateX: translateX.value }],
  }));
  
  return (
    <GestureDetector gesture={panGesture}>
      <Animated.View style={animatedStyle}>
        <Text>Swipe зүүн тийш устгах</Text>
      </Animated.View>
    </GestureDetector>
  );
};
\end{lstlisting}

\textbf{react-native-screens}

React Native Screens нь навигацийн дэлгэцүүдийг native container ашиглан оновчлолдог. Энэ нь санах ойн хэрэглээг багасгаж, дэлгэц хоорондын шилжилтийг хурдасгадаг.

\textbf{react-native-safe-area-context}

Safe Area нь iPhone X-ээс хойших notch (зүсэлт), home indicator (гэрийн товч) зэрэг хэсгүүдэд UI элементүүд давхцахгүй байлгах боломжийг олгодог.

\begin{lstlisting}[language=JavaScript]
import { SafeAreaView, useSafeAreaInsets } from 'react-native-safe-area-context';

// Бүтэн дэлгэц safe area
const Screen = () => (
  <SafeAreaView style={{ flex: 1 }}>
    <Text>Контент</Text>
  </SafeAreaView>
);

// Hook ашиглах
const CustomHeader = () => {
  const insets = useSafeAreaInsets();
  return (
    <View style={{ paddingTop: insets.top }}>
      <Text>Header</Text>
    </View>
  );
};
\end{lstlisting}

\subsubsection{Хөгжүүлэлтийн хэрэгслүүд (Development Tools)}
Кодын чанарыг хангах, алдаа илрүүлэхэд дараах хэрэгслүүдийг ашигласан. Хөгжүүлэлтийн хэрэгслүүд нь кодын чанарыг дээшлүүлж, багийн бүтээмжийг нэмэгдүүлдэг.

\begin{table}[H]
\centering
\caption{Хөгжүүлэлтийн хэрэгслүүд}
\label{tab:dev-tools}
\begin{tabular}{|l|l|p{6cm}|}
\hline
\textbf{Хэрэгсэл} & \textbf{Хувилбар} & \textbf{Зориулалт} \\
\hline
ESLint & \textasciicircum9.25.0 & Код шалгалт (Linting) \\
\hline
eslint-config-expo & \textasciitilde10.0.0 & Expo ESLint тохиргоо \\
\hline
TypeScript & \textasciitilde5.9.2 & Static type checking \\
\hline
@types/react & \textasciitilde19.1.0 & React-ийн төрлийн тодорхойлолт \\
\hline
@react-native-community/cli & latest & React Native CLI \\
\hline
\end{tabular}
\end{table}

\textbf{ESLint - Код шалгалт}

ESLint нь JavaScript/TypeScript кодын алдаа, стилийн асуудлыг автоматаар илрүүлдэг. Энэ нь:
\begin{itemize}
    \item Синтаксийн алдаа олох
    \item Ашиглагдаагүй хувьсагч илрүүлэх
    \item Кодын стандарт мөрдүүлэх
    \item Potential bugs (боломжит алдаа) олох
\end{itemize}

\textbf{ESLint дүрмийн жишээ:}
\begin{lstlisting}[language=JavaScript]
// .eslintrc.js
module.exports = {
  extends: ['expo', 'prettier'],
  plugins: ['react-hooks'],
  rules: {
    // React Hooks дүрмүүд
    'react-hooks/rules-of-hooks': 'error',
    'react-hooks/exhaustive-deps': 'warn',
    
    // Ашиглагдаагүй хувьсагч
    'no-unused-vars': 'warn',
    
    // Console.log хориглох (production-д)
    'no-console': 'warn',
    
    // == биш === ашиглах
    'eqeqeq': 'error',
  },
};
\end{lstlisting}

\textbf{TypeScript - Статик төрөл шалгалт}

TypeScript нь компайл хийх үед төрлийн алдааг илрүүлдэг. Энэ нь runtime алдааг урьдчилан сэргийлж, кодын найдвартай байдлыг нэмэгдүүлдэг.

\textbf{TypeScript тохиргоо (tsconfig.json):}
\begin{lstlisting}[language=JavaScript]
{
  "compilerOptions": {
    "target": "ES2020",
    "module": "ESNext",
    "strict": true,           // Хатуу төрөл шалгалт
    "noImplicitAny": true,    // any төрөл хориглох
    "strictNullChecks": true, // null/undefined шалгах
    "esModuleInterop": true,
    "skipLibCheck": true,
    "forceConsistentCasingInFileNames": true,
    "jsx": "react-native",
    "moduleResolution": "bundler",
    "paths": {
      "@/*": ["./src/*"]      // Path alias
    }
  },
  "include": ["**/*.ts", "**/*.tsx"],
  "exclude": ["node_modules"]
}
\end{lstlisting}

\textbf{TypeScript-ийн жишээ Tdelivery системд:}
\begin{lstlisting}[language=JavaScript]
// Захиалгын төрөл тодорхойлох
interface Order {
  id: string;
  customerId: string;
  items: OrderItem[];
  status: OrderStatus;
  totalAmount: number;
  createdAt: Date;
  deliveryAddress: Address;
}

interface OrderItem {
  productId: string;
  productName: string;
  quantity: number;
  unitPrice: number;
}

type OrderStatus = 'pending' | 'confirmed' | 'delivering' | 'delivered' | 'cancelled';

interface Address {
  street: string;
  district: string;
  city: string;
  coordinates?: {
    latitude: number;
    longitude: number;
  };
}

// Функцийн төрөл
const calculateOrderTotal = (items: OrderItem[]): number => {
  return items.reduce((sum, item) => sum + item.quantity * item.unitPrice, 0);
};

// Generic функц
const getById = async <T>(collection: string, id: string): Promise<T | null> => {
  const doc = await getDoc(doc(db, collection, id));
  return doc.exists() ? (doc.data() as T) : null;
};
\end{lstlisting}

\textbf{VS Code - Хөгжүүлэлтийн орчин}

Visual Studio Code нь Microsoft-оос хөгжүүлсэн үнэгүй код засварлагч бөгөөд React Native хөгжүүлэлтэд хамгийн тохиромжтой. Дараах өргөтгөлүүдийг ашигласан:
\begin{itemize}
    \item \textbf{ESLint:} Кодын алдаа харуулах
    \item \textbf{Prettier:} Кодын формат автоматаар засах
    \item \textbf{React Native Tools:} Debug, log харах
    \item \textbf{GitLens:} Git түүхийг харах
    \item \textbf{Tailwind CSS IntelliSense:} NativeWind class-уудын автомат бөглөлт
\end{itemize>

\textbf{Git - Хувилбарын удирдлага}

Git нь кодын хувилбарыг хянах систем бөгөөд олон хөгжүүлэгч хамтран ажиллахад зайлшгүй шаардлагатай. Tdelivery төсөлд GitHub дээр хадгалагдсан бөгөөд дараах branch стратеги ашигласан:
\begin{itemize}
    \item \textbf{main:} Production-д гарсан код
    \item \textbf{develop:} Хөгжүүлэлтийн үндсэн branch
    \item \textbf{feature/*:} Шинэ боломж нэмэх branch
    \item \textbf{bugfix/*:} Алдаа засах branch
\end{itemize}

\subsubsection{Технологийн стекийн нэгтгэл}
Дээрх бүх технологиуд хамтдаа ажиллаж, бүрэн функциональ B2B түгээлтийн систем болгодог. Технологиудыг зөв сонгох, нэгтгэх нь төслийн амжилтын гол түлхүүр юм.

\textbf{Технологийн сонголтын шалгуурууд:}
\begin{enumerate}
    \item \textbf{Community support:} Идэвхтэй коммьюнити, тогтмол шинэчлэлт
    \item \textbf{Documentation:} Сайн баримт бичиг
    \item \textbf{Performance:} Гүйцэтгэлийн өндөр түвшин
    \item \textbf{Scalability:} Өргөтгөх боломж
    \item \textbf{Developer Experience:} Хөгжүүлэгчийн туршлага
\end{enumerate}

\textbf{Технологиудын харилцан үйлчлэл:}

\begin{figure}[H]
\centering
\begin{tikzpicture}[
    layer/.style={rectangle, draw, minimum width=12cm, minimum height=1.2cm, align=center, font=\small},
    arrow/.style={->, thick}
]
% Layers from bottom to top
\node[layer, fill=blue!20] (backend) at (0,0) {Backend: Firebase (Firestore, Auth, Storage, Functions)};
\node[layer, fill=green!20] (state) at (0,1.5) {State Management: React Context API + AsyncStorage};
\node[layer, fill=yellow!20] (logic) at (0,3) {Business Logic: TypeScript Services};
\node[layer, fill=orange!20] (ui) at (0,4.5) {UI Layer: React Native + NativeWind + React Native Paper};
\node[layer, fill=red!20] (platform) at (0,6) {Platform: Expo SDK (iOS + Android + Web)};
\end{tikzpicture}
\caption{Технологийн стекийн архитектур}
\label{fig:tech-stack}
\end{figure}

\textbf{Давхаргуудын тайлбар:}

\begin{itemize}
    \item \textbf{Platform Layer (Платформын давхарга):} Expo SDK нь iOS, Android, Web гурван платформд нэг кодоор ажиллах боломжийг олгоно. Expo-ийн бэлэн модулиуд (Camera, Location, Notifications) нь native код бичихгүйгээр төхөөрөмжийн боломжуудад хандана.
    
    \item \textbf{UI Layer (Хэрэглэгчийн интерфэйсийн давхарга):} React Native компонентууд, React Native Paper (Material Design), NativeWind (TailwindCSS) хамтран ажиллаж гоёмсог, responsive UI үүсгэнэ. Reanimated болон Gesture Handler нь анимаци, жестүүдийг удирдана.
    
    \item \textbf{Business Logic Layer (Бизнес логикийн давхарга):} TypeScript-ээр бичигдсэн service-үүд нь бизнес дүрмүүдийг хэрэгжүүлнэ. Жишээ нь: захиалгын тооцоо, хүргэлтийн төлөв шинэчлэлт, мэдэгдэл илгээх логик.
    
    \item \textbf{State Management Layer (Төлөв удирдлагын давхарга):} React Context API нь global state (хэрэглэгчийн мэдээлэл, сагс, тохиргоо) удирдана. AsyncStorage нь offline горимд өгөгдөл хадгална.
    
    \item \textbf{Backend Layer (Backend давхарга):} Firebase үйлчилгээнүүд (Firestore, Authentication, Storage, Cloud Functions) нь серверийн бүх функцийг хариуцна. Real-time синхрончлол, offline support, аюулгүй байдал бүгд Firebase-аар хангагдана.
\end{itemize}

\textbf{Өгөгдлийн урсгал (Data Flow):}

\begin{enumerate}
    \item Хэрэглэгч UI дээр үйлдэл хийнэ (товч дарах, form бөглөх)
    \item UI компонент event handler-ийг дуудна
    \item Handler нь Context-ийн функцийг дуудна
    \item Context нь Firebase SDK-аар серверт хүсэлт илгээнэ
    \item Firebase өгөгдлийг боловсруулж, хариу буцаана
    \item Context state-ийг шинэчилнэ
    \item React компонентуудыг дахин render хийнэ
    \item UI шинэчлэгдэнэ
\end{enumerate}

\textbf{Технологийн нэгтгэлийн үр дүн:}

\begin{table}[H]
\centering
\caption{Технологийн нэгтгэлийн үр дүн}
\label{tab:tech-benefits}
\begin{tabular}{|l|p{10cm}|}
\hline
\textbf{Үзүүлэлт} & \textbf{Үр дүн} \\
\hline
Хөгжүүлэлтийн хугацаа & Нэг платформд хөгжүүлэхтэй харьцуулахад 40-50\% бага \\
\hline
Кодын дахин ашиглалт & iOS, Android, Web хооронд 85-90\% код дахин ашиглагдсан \\
\hline
Гүйцэтгэл & 60 FPS, native аппликейшнтэй дүйцэхүйц \\
\hline
Offline support & Интернэтгүй үед ч үндсэн функцууд ажиллана \\
\hline
Real-time & Захиалгын төлөв өөрчлөгдөхөд 100ms дотор мэдэгдэл хүрнэ \\
\hline
\end{tabular}
\end{table}

\subsection{Ижил төстэй системүүдийн харьцуулалт}
Дэлхийн зах зээлд болон бүс нутгийн хэмжээнд ашиглагдаж буй ижил төстэй системүүдийг харьцуулан судлав.

\begin{longtable}{|p{2.5cm}|p{3.5cm}|p{3.5cm}|p{3.5cm}|}
\caption{Ижил төстэй системүүдийн харьцуулалт} \label{tab:comparison} \\
\hline
\textbf{Шалгуур} & \textbf{Lalamove (Ази)} & \textbf{Borzo (ОХУ/Монгол)} & \textbf{Tdelivery (Санал болгож буй)} \\
\hline
\endfirsthead
\hline
\textbf{Шалгуур} & \textbf{Lalamove (Ази)} & \textbf{Borzo (ОХУ/Монгол)} & \textbf{Tdelivery (Санал болгож буй)} \\
\hline
\endhead
\hline
\endfoot
\hline
\endlastfoot

\textbf{Үндсэн зорилго} & C2C, B2C Шуурхай хүргэлт & Crowdsourced хүргэлт & B2B Бараа түгээлт, SCM \\
\hline
\textbf{Хэрэглэгчид} & Хувь хүн, Жижиг бизнес & Хувь хүн, Онлайн дэлгүүр & Бөөний төв, Дэлгүүр, Нийлүүлэгч \\
\hline
\textbf{Тээврийн хэрэгсэл} & Мотоцикл, Ван, Машин & Явган, Мото, Машин & Портер, Ван (Ачаа тээвэр) \\
\hline
\textbf{Real-time Tracking} & Маш сайн & Сайн & Сайн \\
\hline
\textbf{Маршрут оновчлол} & AI суурьтай, автомат хуваарилалт & Автомат хуваарилалт & Google Maps API, Гараар болон автомат \\
\hline
\textbf{Төлбөрийн шийдэл} & Онлайн, Бэлэн & Бэлэн, Данс, COD & Бэлэн, Данс, Зээл (Credit) \\
\hline
\textbf{e-POD} & Зураг, Гарын үсэг & Зураг & Зураг, Гарын үсэг, GPS \\
\hline
\textbf{Үнэ} & Зайнаас хамаарсан & Зайнаас хамаарсан & Захиалгын дүнгийн хувь, эсвэл сарын хураамж \\
\hline
\textbf{Сул тал} & B2B нарийн процесс (нэхэмжлэх, буцаалт) дутмаг & Тогтмол хүргэгчгүй, найдвартай байдал сул & Зах зээлд шинэ, хэрэглэгчийн бааз бага \\
\hline
\end{longtable}

\subsubsection{Харьцуулсан дүгнэлт}
Lalamove болон Borzo нь "On-demand" буюу яаралтай, нэг удаагийн хүргэлтэд маш сайн шийдэл боловч B2B харилцааны нарийн процессуудыг (бараа буцаалт, зээлийн тооцоо, тогтмол маршрут, их хэмжээний бараа хүлээлцэх) бүрэн шийдэж чадахгүй байна. "Tdelivery" систем нь эдгээр дутагдлыг нөхөж, Монголын жижиг дунд бизнесүүдийн онцлогт тохирсон, нийлүүлэлтийн сүлжээг бүхэлд нь хамарсан шийдэл болохоороо давуу талтай юм.

\vfill
%-------------------------------------------------------------------------------
%   CHAPTER 3 - METHODOLOGY
%-------------------------------------------------------------------------------

\section{Судалгааны арга зүй}

\subsection{Хөгжүүлэлтийн арга зүй}
Програм хангамжийг хөгжүүлэхдээ **Agile** арга зүйг, тэр дундаа **Scrum** фреймворкийг ашигласан. Энэ нь төслийг бага хэмжээний, удирдах боломжтой хэсгүүдэд (Sprints) хувааж, тогтмол хугацаанд үр дүнг гаргах, хэрэглэгчийн санал хүсэлтэд тулгуурлан сайжруулалт хийх боломжийг олгодог.

\subsection{Системийн шинжилгээ}
Системийн шаардлагыг тодорхойлохдоо хэрэглэгчийн түүх (User Stories) болон хэрэглээний тохиолдол (Use Cases) аргуудыг ашигласан.

\subsubsection{Хэрэглэгчийн ангилал}
\begin{itemize}
    \item \textbf{Захиалагч (Customer):} Бараа захиалах дэлгүүр, ресторан, хувь хүн.
    \item \textbf{Нийлүүлэгч (Supplier):} Бараа нийлүүлэгч компани, бөөний төв.
    \item \textbf{Хүргэгч (Courier):} Барааг хүргэх үүрэгтэй ажилтан.
    \item \textbf{Админ (Admin):} Системийн хэвийн үйл ажиллагааг хангах, хэрэглэгчдийг удирдах.
\end{itemize}

\subsubsection{Функциональ шаардлага}
Систем нь дараах үндсэн функциональ шаардлагуудыг хангасан байна.

\begin{longtable}{|p{1.5cm}|p{4cm}|p{7cm}|}
\caption{Функциональ шаардлагууд} \label{tab:func_req} \\
\hline
\textbf{ID} & \textbf{Шаардлага} & \textbf{Тайлбар} \\
\hline
\endfirsthead
\hline
\textbf{ID} & \textbf{Шаардлага} & \textbf{Тайлбар} \\
\hline
\endhead
\hline
\endfoot
\hline
\endlastfoot

FR-01 & Бүртгэл ба Нэвтрэх & Хэрэглэгч утасны дугаар болон имэйлээр бүртгүүлэх, нэвтрэх боломжтой байх. \\
\hline
FR-02 & Бараа хайх & Захиалагч барааг нэр, төрөл, үнээр шүүж хайх. \\
\hline
FR-03 & Сагс удирдах & Барааг сагсанд нэмэх, тоо хэмжээг өөрчлөх, хасах. \\
\hline
FR-04 & Захиалга үүсгэх & Сагсан дахь барааг баталгаажуулж захиалга илгээх. \\
\hline
FR-05 & Захиалга хянах & Захиалгын төлөвийг (Pending, Confirmed, Delivering, Delivered) бодит цагт харах. \\
\hline
FR-06 & Захиалга удирдах & Нийлүүлэгч ирсэн захиалгыг хүлээн авах эсвэл цуцлах. \\
\hline
FR-07 & Хүргэлт хуваарилах & Нийлүүлэгч захиалгыг хүргэгчид хуваарилах. \\
\hline
FR-08 & e-POD бүртгэх & Хүргэгч барааг хүлээлгэн өгөхдөө зураг авч, гарын үсэг зуруулж баталгаажуулах. \\
\hline
\end{longtable}

\subsubsection{Функциональ бус шаардлага}
\begin{itemize}
    \item \textbf{Найдвартай байдал:} Систем нь 99.9\%-ийн ажиллах хугацаатай (uptime) байна.
    \item \textbf{Аюулгүй байдал:} Хэрэглэгчийн нууц үг шифрлэгдсэн байх, өгөгдлийн хандалт нь дүрэмд (Security Rules) захирагдах.
    \item \textbf{Гүйцэтгэл:} Үйлдлийн хариу өгөх хугацаа 2 секундээс хэтрэхгүй байх.
    \item \textbf{Өргөтгөх боломж:} Хэрэглэгчийн тоо нэмэгдэхэд систем гацахгүй ажиллах (Scalability).
\end{itemize}

\subsubsection{Хэрэглээний тохиолдлын диаграмм (Use Case Diagram)}
Системийн үндсэн хэрэглэгчид болон тэдгээрийн үйлдэл, харилцан хамаарлыг дараах диаграммаар харуулав.

\begin{figure}[h]
    \centering
    \begin{tikzpicture}[scale=0.8, every node/.style={scale=0.8}]
        % Actors
        \node[draw, circle, minimum size=1cm] (customer) at (0, 4) {Customer};
        \node[draw, circle, minimum size=1cm] (supplier) at (10, 4) {Supplier};
        \node[draw, circle, minimum size=1cm] (courier) at (5, -2) {Courier};

        % Use Cases
        \node[draw, ellipse] (login) at (5, 6) {Login/Register};
        \node[draw, ellipse] (order) at (2, 2) {Place Order};
        \node[draw, ellipse] (manage) at (8, 2) {Manage Products};
        \node[draw, ellipse] (accept) at (8, 0) {Accept Order};
        \node[draw, ellipse] (deliver) at (5, 0) {Deliver Order};
        \node[draw, ellipse] (track) at (2, 0) {Track Order};

        % Connections
        \draw (customer) -- (login);
        \draw (supplier) -- (login);
        \draw (courier) -- (login);
        
        \draw (customer) -- (order);
        \draw (customer) -- (track);
        
        \draw (supplier) -- (manage);
        \draw (supplier) -- (accept);
        
        \draw (courier) -- (deliver);
        \draw (courier) -- (track);
    \end{tikzpicture}
    \caption{Системийн Use Case диаграмм}
    \label{fig:usecase}
\end{figure}

\subsubsection{Үйл ажиллагааны диаграмм (Activity Diagram)}
Захиалга үүсгэхээс эхлээд хүргэгдэх хүртэлх үйл явцыг Activity диаграммаар дүрслэв.

\begin{figure}[h]
    \centering
    \begin{tikzpicture}[node distance=1.5cm, auto]
        % Nodes
        \node[draw, rounded corners] (start) {Start};
        \node[draw, rectangle, below of=start] (place) {Place Order};
        \node[draw, diamond, below of=place, aspect=2] (check) {Stock Available?};
        \node[draw, rectangle, left of=check, node distance=3cm] (cancel) {Cancel Order};
        \node[draw, rectangle, below of=check, node distance=2cm] (notify) {Notify Supplier};
        \node[draw, rectangle, below of=notify] (prepare) {Prepare Order};
        \node[draw, rectangle, below of=prepare] (assign) {Assign Courier};
        \node[draw, rectangle, below of=assign] (deliver) {Deliver to Customer};
        \node[draw, rounded corners, below of=deliver] (end) {End};

        % Edges
        \draw[->] (start) -- (place);
        \draw[->] (place) -- (check);
        \draw[->] (check) -- node {No} (cancel);
        \draw[->] (check) -- node {Yes} (notify);
        \draw[->] (notify) -- (prepare);
        \draw[->] (prepare) -- (assign);
        \draw[->] (assign) -- (deliver);
        \draw[->] (deliver) -- (end);
        \draw[->] (cancel) |- (end);
    \end{tikzpicture}
    \caption{Захиалгын үйл явцын Activity диаграмм}
    \label{fig:activity}
\end{figure}

\subsection{Системийн архитектур}
"Tdelivery" систем нь орчин үеийн **Serverless** архитектурт суурилсан. Энэ нь сервер талын дэд бүтцийг удирдах шаардлагагүйгээр, зөвхөн бизнесийн логик болон код бичихэд төвлөрөх боломжийг олгодог.

\begin{figure}[h]
    \centering
    \begin{tikzpicture}[node distance=2cm]
        \node[draw, rectangle, minimum width=2.5cm, minimum height=1cm] (app) {Mobile App (React Native)};
        \node[draw, rectangle, right of=app, node distance=4cm] (auth) {Firebase Auth};
        \node[draw, rectangle, below of=auth] (firestore) {Cloud Firestore};
        \node[draw, rectangle, below of=firestore] (functions) {Cloud Functions};
        \node[draw, cylinder, shape border rotate=90, draw, minimum height=1cm, minimum width=1.5cm, right of=firestore, node distance=4cm] (db) {NoSQL DB};

        \draw[<->] (app) -- (auth);
        \draw[<->] (app) -- (firestore);
        \draw[<->] (firestore) -- (functions);
        \draw[<->] (firestore) -- (db);
    \end{tikzpicture}
    \caption{Системийн архитектур}
    \label{fig:architecture}
\end{figure}

\subsubsection{Технологийн сонголт}
\begin{itemize}
    \item \textbf{Frontend (Mobile App):} React Native (Expo SDK). Кросс-платформ хөгжүүлэлт хийхэд тохиромжтой, хурдан, өндөр гүйцэтгэлтэй.
    \item \textbf{Backend (BaaS):} Google Firebase.
        \begin{itemize}
            \item \textbf{Authentication:} Хэрэглэгчийн бүртгэл, нэвтрэх эрх.
            \item \textbf{Firestore:} NoSQL өгөгдлийн сан.
            \item \textbf{Cloud Functions:} Сервер талын логик.
            \item \textbf{Storage:} Зураг, файл хадгалах.
        \end{itemize}
\end{itemize}

\subsection{Өгөгдлийн сангийн зохиомж (ERD)}
Систем нь NoSQL төрлийн Cloud Firestore өгөгдлийн санг ашиглана. Өгөгдлийн хамаарлыг ER диаграммаар харуулав.

\begin{figure}[h]
    \centering
    \begin{tikzpicture}[node distance=2cm]
        \node[draw, rectangle] (user) {User};
        \node[draw, rectangle, right of=user, node distance=4cm] (order) {Order};
        \node[draw, rectangle, below of=order] (product) {Product};
        
        \draw[->] (user) -- node[above] {places} (order);
        \draw[->] (user) -- node[below] {supplies} (product);
        \draw[->] (order) -- node[right] {contains} (product);
    \end{tikzpicture}
    \caption{Өгөгдлийн сангийн ER диаграмм (Хялбаршуулсан)}
    \label{fig:erd}
\end{figure}

\subsubsection{Өгөгдлийн сангийн бүтэц (Schema)}
Firestore нь схемгүй (schemaless) боловч өгөгдлийн бүрэн бүтэн байдлыг хангахын тулд дараах бүтцийг баримтална.

\textbf{1. Users Collection}
\begin{itemize}
    \item \texttt{uid} (String): Хэрэглэгчийн ID (PK)
    \item \texttt{email} (String): Имэйл хаяг
    \item \texttt{role} (String): 'customer' | 'supplier' | 'courier'
    \item \texttt{displayName} (String): Нэр
    \item \texttt{phoneNumber} (String): Утасны дугаар
    \item \texttt{address} (Map): Хаягийн мэдээлэл (lat, lng, street)
\end{itemize}

\textbf{2. Products Collection}
\begin{itemize}
    \item \texttt{id} (String): Бүтээгдэхүүний ID (PK)
    \item \texttt{supplierId} (String): Нийлүүлэгчийн ID (FK)
    \item \texttt{name} (String): Нэр
    \item \texttt{price} (Number): Үнэ
    \item \texttt{stock} (Number): Үлдэгдэл
    \item \texttt{imageUrl} (String): Зургийн зам
    \item \texttt{category} (String): Ангилал
\end{itemize}

\textbf{3. Orders Collection}
\begin{itemize}
    \item \texttt{id} (String): Захиалгын ID (PK)
    \item \texttt{customerId} (String): Захиалагчийн ID (FK)
    \item \texttt{supplierId} (String): Нийлүүлэгчийн ID (FK)
    \item \texttt{items} (Array): Захиалсан бараанууд [{productId, quantity, price}]
    \item \texttt{totalAmount} (Number): Нийт дүн
    \item \texttt{status} (String): 'pending' | 'confirmed' | 'delivering' | 'delivered' | 'cancelled'
    \item \texttt{createdAt} (Timestamp): Үүсгэсэн огноо
\end{itemize}

\subsection{Алгоритм ба Логик шийдэл}
Системийн үр ашигтай ажиллагааг хангахын тулд хэд хэдэн алгоритм, логик шийдлүүдийг хэрэгжүүлсэн.

\subsubsection{Атомик Денормализаци (Atomic Denormalization)}
NoSQL өгөгдлийн санд өгөгдлийг унших хурдыг нэмэгдүүлэхийн тулд денормализаци хийх нь түгээмэл байдаг. Гэвч өгөгдлийн зөрчил үүсэх эрсдэлтэй. Үүнийг шийдэхийн тулд Firestore Transaction ашиглан "Атомик Денормализаци" аргыг хэрэгжүүлсэн.

\subsubsection{Бодит цагийн захиалгын удирдлага}
Захиалгын төлөв өөрчлөгдөх бүрт хэрэглэгчдэд шууд харагдах ёстой. Үүний тулд \texttt{onSnapshot} сонсогч (listener) ашигласан.

\subsection{Тестлэлт}
Системийн чанарыг баталгаажуулахын тулд Unit Testing, Integration Testing, User Acceptance Testing (UAT) зэрэг аргуудыг ашигласан.

\vfill

%-------------------------------------------------------------------------------
%   CHAPTER 4 - EXPERIMENTS AND RESULTS (SHORTENED)
%-------------------------------------------------------------------------------

\section{Туршилт, үр дүн}
\label{sec:experiments}

Тус бүлэгт Бүлэг 1-д дэвшүүлсэн RQ1--RQ3 асуултад хариулах туршилтын үр дүнг танилцуулна: (i) FCFS claim загвар \texttt{FOR UPDATE SKIP LOCKED}-оор race condition-ыг шийдэж буй эсэх (RQ1); (ii) RLS нь multi-tenant isolation-ыг application-layer алдаанаас үл хамааран хангаж буй эсэх (RQ2); (iii) системийн performance (p95 < 2~с, concurrency safe) шаардлагад нийцэж буй эсэх (RQ3-ийн дэд асуулт). Туршилтыг concurrency, load, penetration, functional (E2E) гэсэн дөрөвт хуваав.

%===============================================================================
\subsection{Туршилтын орчин ба аргачлал}
\label{subsec:test_environment}

Туршилтыг production-аас тусгаарласан staging орчинд явуулсан. Орчны үзүүлэлт: Supabase Pro tier, PostgreSQL 15.6, PgBouncer transaction pooling, pool size 15. Клиент машин: MacBook Pro (Apple M1, 16~GB RAM), Node.js 20.11. Клиент — Supabase дундаж RTT 35--45~ms. Production schema, RLS бодлого, RPC функц, триггерүүдийг staging-т хуулбарласан.

Туршилтын өгөгдлийг \texttt{@faker-js/faker}-аар үүсгэв: 20 байгууллага, 200 хэрэглэгч (100 борлуулагч + 100 хүргэгч), 500 бүтээгдэхүүн, 1{,}000 захиалга, 1{,}000 даалгавар. Хүргэгчдийн 80\%-ийг \texttt{approved}, 20\%-ийг \texttt{pending} төлөвт санаатай үлдээсэн нь RLS-ийн approved-only шүүлтийг шалгах зорилготой. Туршилт бүрийн өмнө \texttt{TRUNCATE ... RESTART IDENTITY CASCADE}-ээр өгөгдлийг анхны төлөвт буцаасан.

Ашигласан хэрэгсэл: \texttt{@supabase/supabase-js 2.46} (concurrency test, \texttt{Promise.all()}), \texttt{k6 0.54}~\cite{k6docs} (load test, p50/p95/p99 хэмжилт), \texttt{psql} + \texttt{SET request.jwt.claims} (RLS verification), захиалгат JS script-ууд (penetration, OTP brute-force, E2E). Бүх script, seed файл, k6 тест Git-ийн \texttt{tests/} санд байршиж давтагдах боломжтой.

%===============================================================================
\subsection{Зэрэгцээ хандалтын туршилт}
\label{subsec:concurrency_test}

Зорилго: $N$ хүргэгч нэг даалгаврыг зэрэг claim хийхэд зөвхөн нэг нь амжилттай болж, үлдсэн нь чимээгүйгээр татгалзагдах эсэхийг шалгах. Арга зүй: (1)~шинэ \texttt{delivery\_tasks} мөр \texttt{status='open'}-оор үүсгэх, (2)~$N$ хүргэгчийн JWT токеноор client instance бэлтгэх, (3)~\texttt{Promise.all()}-оор $N$ RPC дуудлагыг зэрэг илгээх, (4)~өгөгдлийн сангаас \texttt{SELECT courier\_id}-г шалгаж нэг мөр л бичигдсэн invariant-ыг батлах. $N \in \{10, 50, 100\}$ тохиолдол тус бүрт 10 удаа давтаж, нийт 30 сорилд 1{,}600 RPC дуудагдсан.

\begin{table}[H]
\centering
\renewcommand{\arraystretch}{1.25}
\caption{Зэрэгцээ claim туршилтын үр дүн ($N$ хүргэгч, 10 давталт)}
\label{tab:concurrency_results}
\footnotesize
\begin{tabular}{|c|c|c|c|c|c|}
\hline
$N$ & Давталт & Амжилтын дундаж & Давхардсан & Амжилтын \% & Дундаж хугацаа \\
\hline
10  & 10 & 1.00 & 0 & 100\% & 124~ms \\
\hline
50  & 10 & 1.00 & 0 & 100\% & 287~ms \\
\hline
100 & 10 & 1.00 & 0 & 100\% & 512~ms \\
\hline
\end{tabular}
\end{table}

Нийт 30 сорилын туршид \textbf{давхардсан claim нэг ч удаа гараагүй}. Үлдсэн $N{-}1$ хүсэлт \texttt{NULL} буцаасан. Клиент тал үүнийг ``даалгавар авагдсан'' мэдэгдэл болгон хөрвүүлнэ. 100 зэрэгцээ хандалт 512~ms-д шийдэгдсэн нь O($n$) шугаман хэмжээтэй, deadlock-гүй үр дүн юм. Энэ нь pessimistic locking-ийн онолын хүлээгдсэн зан төлөвтэй нийцнэ~\cite{bernstein1987concurrency, kleppmann2017ddia}.

Хэрэв claim нь ердийн \texttt{UPDATE ... WHERE status='open'}-оор хэрэгжсэн бол MVCC-ийн улмаас хоёр зэрэгцээ UPDATE хоёулаа мөрийг ``нээлттэй'' гэж үзэж \textit{lost update} үүсэх магадлалтай. Иймд \texttt{FOR UPDATE SKIP LOCKED} нь зөвхөн performance-ийн биш, \textit{correctness}-ийн хувьд зайлшгүй шийдэл юм.

%===============================================================================
\subsection{Гүйцэтгэлийн туршилт (Load Test)}
\label{subsec:load_test}

Гүйцэтгэлийн туршилтыг k6-аар 100 VU, 10~мин ramp-up + 20~мин тогтмол ачаалалтайгаар ажиллуулсан. Үр дүнг Хүснэгт~\ref{tab:endpoint_latency}-д нэгтгэв.

\begin{table}[H]
\centering
\renewcommand{\arraystretch}{1.25}
\caption{k6 ачааллын туршилтын endpoint-үүдийн хариу хугацаа (ms)}
\label{tab:endpoint_latency}
\footnotesize
\begin{tabular}{|p{4.8cm}|c|c|c|c|c|}
\hline
\textbf{Endpoint} & \textbf{p50} & \textbf{p95} & \textbf{p99} & \textbf{Max} & \textbf{Амжилт} \\
\hline
\texttt{GET available\_tasks}      & 145 & 312 & 456 & 892   & 100\% \\
\hline
\texttt{POST claim\_delivery\_task} & 287 & 621 & 834 & 1\,200 & 100\% \\
\hline
\texttt{POST verify\_epod\_otp}    & 198 & 287 & 345 & 523   & 100\% \\
\hline
\texttt{GET orders} (борлуулагч)   & 167 & 389 & 512 & 734   & 100\% \\
\hline
Dashboard KPI агрегат              & 234 & 567 & 789 & 1\,100 & 100\% \\
\hline
\end{tabular}
\end{table}

Хүснэгтийн үр дүнг визуал хэлбэрээр Зураг~\ref{fig:endpoint_latency_chart}-д харуулав.

\begin{figure}[H]
\centering
\resizebox{0.85\textwidth}{!}{%
\begin{tikzpicture}
\begin{axis}[
    ybar,
    bar width=12pt,
    width=0.95\textwidth,
    height=7.5cm,
    enlarge x limits=0.12,
    ymin=0, ymax=900,
    ylabel={Хариу хугацаа (ms)},
    symbolic x coords={available\_tasks,claim\_task,verify\_otp,orders,dashboard},
    xtick=data,
    x tick label style={font=\small, rotate=15, anchor=east},
    ymajorgrids=true,
    grid style={gray!30},
    legend style={at={(0.5,-0.25)}, anchor=north, legend columns=3, font=\small},
    ylabel style={font=\small},
    tick label style={font=\small}
]
\addplot[fill=blue!40, draw=blue!60!black] coordinates {
    (available\_tasks,145) (claim\_task,287) (verify\_otp,198) (orders,167) (dashboard,234)
};
\addplot[fill=orange!60, draw=orange!70!black] coordinates {
    (available\_tasks,312) (claim\_task,621) (verify\_otp,287) (orders,389) (dashboard,567)
};
\addplot[fill=red!50, draw=red!70!black] coordinates {
    (available\_tasks,456) (claim\_task,834) (verify\_otp,345) (orders,512) (dashboard,789)
};
\legend{p50, p95, p99}
\end{axis}
\end{tikzpicture}
}
\caption{k6 load test — 5 endpoint-ийн p50/p95/p99 latency хуваарилалт (100 VU, 30 мин)}
\label{fig:endpoint_latency_chart}
\end{figure}

Гол ажиглалтууд:
\begin{itemize}
  \item \texttt{available\_tasks} нь \texttt{(status)} талбарын partial index ашигладаг тул p95 312~ms-д багтсан.
  \item \texttt{claim\_delivery\_task}-ийн p95 621~ms нь lock acquisition-оос ихээр хамаарч, \S\ref{subsec:concurrency_test}-ийн үр дүнтэй нийцнэ.
  \item \texttt{verify\_epod\_otp}-ийн p50 198~ms-ын ойролцоогоор 160~ms нь bcrypt (cost 10) hash comparison-д зарцуулагдана. Энэ нь brute-force эсэргүүцэхэд зориудаар тавьсан зардал юм~\cite{provos1999bcrypt}.
  \item Бүх endpoint \textbf{p95 < 2~с} (хатуу хязгаар), \textbf{p99 < 1.5~с} (зөөлөн хязгаар)-ыг хангасан. \texttt{claim}-ийн max 1.2~с нь edge case-ийг анхаарах шаардлагатайг илтгэнэ.
\end{itemize}
Иймд 100 VU ачаалал дор performance шаардлага биелсэн гэж дүгнэв.

%===============================================================================
\subsection{Аюулгүй байдлын туршилт}
\label{subsec:security_test}

\subsubsection{RLS бодлогын нэвтрэлтийн тест}

Org~A хэрэглэгчийн JWT-ээр Org~B-гийн өгөгдөлд хандах 6 хувилбарыг туршив. Хүлээгдэж буй үр дүн: бүх тохиолдолд RLS хандалтыг хаах, эсвэл 0 мөр буцаах.

\begin{table}[H]
\centering
\renewcommand{\arraystretch}{1.25}
\caption{RLS бодлогын нэвтрэлтийн тестийн үр дүн}
\label{tab:rls_pentest}
\footnotesize
\begin{tabular}{|c|p{7cm}|p{4cm}|c|}
\hline
\textbf{ID} & \textbf{Хувилбар} & \textbf{Бодит үр дүн} & \textbf{Төлөв} \\
\hline
T1 & Org~A → Org~B \texttt{orders} SELECT & 0 мөр & $\checkmark$ \\
\hline
T2 & Org~A → Org~B \texttt{delivery\_tasks} UPDATE & 0 мөр шинэчлэгдсэн & $\checkmark$ \\
\hline
T3 & \texttt{pending} хүргэгч \texttt{available\_tasks} SELECT & 0 мөр & $\checkmark$ \\
\hline
T4 & Хүргэгч~A → Хүргэгч~B \texttt{courier\_earnings} SELECT & 0 мөр & $\checkmark$ \\
\hline
T5 & SQL injection (\texttt{.eq('org\_id', "' OR 1=1 --")}) & Parameterized, 0 мөр & $\checkmark$ \\
\hline
T6 & Алдаатай / хуурамч JWT & 401 Unauthorized & $\checkmark$ \\
\hline
\end{tabular}
\end{table}

Бүх 6 тест тэнцсэн. T5 нь PostgREST-ийн prepared statement хандлагыг, T3 нь approved шалгалт клиент биш ӨС талд хийгддэгийг тус тус баталсан.

\subsubsection{OTP brute-force ба JWT баталгаажуулалт}

Идэвхтэй OTP сесст буруу кодыг дараалан 6 удаа илгээв. Тав дахь оролдлогын дараа өгөгдлийн сангийн түвшинд блоклогдсон. Зургаа дахь хүсэлт нь bcrypt comparison хүртэл хүрэлгүй 43~ms-д эрт таслагдсан (1--5 оролдлого: 72--86~ms; 6-р оролдлого: ``Хязгаар хэтэрсэн'' алдаа). Brute-force амжилтын онолын магадлал $P = 5/10^6 = 5 \times 10^{-4}\%$. OTP-ийн 10~мин expiry-тэй хослуулахад практикт мэдэгдэхүйц бус~\cite{owasp2024authcheatsheet}. Expiry шалгах нэмэлт тестэд 11~минутын дараа зөв кодоор илгээсэн хүсэлт ``expired'' алдаагаар татгалзагдсан нь \texttt{expires\_at} шалгалтыг баталсан.

JWT-ийн 4 хэлбэрийн дайралтыг шалгав:
\begin{enumerate}
  \item \texttt{org\_id} claim-ийг өөрчилсөн;
  \item \texttt{role='service\_role'} болгож privilege escalation оролдсон;
  \item \texttt{alg=none} signature-гүй токен;
  \item хугацаа дууссан токен.
\end{enumerate}
Бүгд PostgREST түвшинд 401 Unauthorized-аар татгалзагдсан. Өгөгдлийн сан хүртэл хүсэлт дамжаагүй~\cite{jones2015jwt}.

%===============================================================================
\subsection{End-to-End (E2E) туршилт}
\label{subsec:functional_test}

\S\ref{subsec:system_overview}-ийн 8 алхамт үндсэн урсгалыг end-to-end script-ээр туршиж, цагийн тэмдгүүдийн invariant (\texttt{assigned\_at} < \texttt{picked\_up\_at} < \texttt{delivered\_at} < \texttt{completed\_at}) зөрөлгүй биелсэн. Үр дүнг Хүснэгт~\ref{tab:e2e_flow}-д харуулав.

\begin{table}[H]
\centering
\renewcommand{\arraystretch}{1.2}
\caption{E2E захиалгын урсгалын баталгаажуулалт}
\label{tab:e2e_flow}
\footnotesize
\begin{tabular}{|c|p{4.5cm}|p{5.5cm}|c|}
\hline
\textbf{\#} & \textbf{Алхам} & \textbf{Шалгуур} & \textbf{Төлөв} \\
\hline
1 & Байгууллага бүртгэл & \texttt{orgs} мөр + JWT \texttt{org\_id} claim & $\checkmark$ \\
\hline
2 & Бүтээгдэхүүн нэмэх & \texttt{products} + RLS шүүлт & $\checkmark$ \\
\hline
3 & Захиалга хүлээн авах & \texttt{orders.status=paid} & $\checkmark$ \\
\hline
4 & Даалгавар нийтлэх & \texttt{delivery\_tasks.status=open} & $\checkmark$ \\
\hline
5 & Хүргэгч авах & \texttt{status=assigned}, \texttt{assigned\_at} & $\checkmark$ \\
\hline
6 & Бараа авах & \texttt{status=picked\_up}, \texttt{picked\_up\_at} & $\checkmark$ \\
\hline
7 & Хүргэх & \texttt{status=delivered}, \texttt{delivered\_at} & $\checkmark$ \\
\hline
8 & ePOD + орлого & \texttt{status=completed} + \texttt{courier\_earnings} мөр & $\checkmark$ \\
\hline
\end{tabular}
\end{table}

Серверийн бодит тооцооллын хугацаа 2.4~секунд, нийт урсгал (клиент sleep-тэй) 6.8~секунд байлаа. ePOD дэд урсгал нь 5 алхмыг зөв дарааллаар гүйцэтгэв: OTP үүсгэх → bcrypt hash хадгалах → \texttt{verify\_epod\_otp} RPC → completed trigger → \texttt{record\_courier\_earning} trigger. Энэ нь \S\ref{subsec:security_test}-ийн brute-force хамгаалалттай нэгдмэл ажиллаж буйг баталсан.

%===============================================================================
\subsection{Үр дүнгийн дүгнэлт}
\label{subsec:results_summary}

\subsubsection{Хязгаарлалт}

Туршилт нь (i) Supabase Pro tier-ийн нөөцийн хязгаарт, (ii) synthetic өгөгдөл дээр (peak cyclicity гэх мэт бодит зан төлөвгүй), (iii) single-region тохиргоонд, (iv) 100 VU хүртэлх ачаалалд, (v) тогтвортой сүлжээнд (RTT 35--45~ms) хийгдсэн. 500+ VU, multi-region, 3G/4G сүлжээний нөхцөл үнэлэгдээгүй тул эдгээр нь ирээдүйн өргөтгөлийн судалгааны чиглэл болно.

\subsubsection{Шаардлагын биелэлтийн нэгтгэл}

\begin{table}[H]
\centering
\renewcommand{\arraystretch}{1.25}
\caption{Системийн чанарын шаардлагуудын туршилтаар баталгаажсан биелэлт}
\label{tab:requirement_verification}
\footnotesize
\begin{tabular}{|p{2.5cm}|p{5.5cm}|p{4cm}|c|}
\hline
\textbf{Ангилал} & \textbf{Шаардлага} & \textbf{Баталгаа} & \textbf{Төлөв} \\
\hline
Гүйцэтгэл & available\_tasks < 2~с & Хүснэгт~\ref{tab:endpoint_latency} (p95=312~ms) & $\checkmark$ \\
\hline
Гүйцэтгэл & claim < 1~с & Хүснэгт~\ref{tab:endpoint_latency} (p99=834~ms) & $\checkmark$ \\
\hline
Аюулгүй байдал & JWT validation & \S\ref{subsec:security_test} (4 хувилбар) & $\checkmark$ \\
\hline
Аюулгүй байдал & RLS тусгаарлалт & Хүснэгт~\ref{tab:rls_pentest} (T1--T5) & $\checkmark$ \\
\hline
Аюулгүй байдал & bcrypt + HTTPS & OTP тулгалт + TLS 1.3 & $\checkmark$ \\
\hline
Өргөтгөх боломж & Зэрэгцээ claim хамгаалалт & Хүснэгт~\ref{tab:concurrency_results} (100 VU, 0 давхардал) & $\checkmark$ \\
\hline
Өргөтгөх боломж & Олон байгууллага зэрэг & T1, T4 + E2E & $\checkmark$ \\
\hline
Хүртээмж & Сессийн автомат сэргээлт & E2E: токен refresh зөв & $\checkmark$ \\
\hline
Ашиглах чадвар & Хүргэгч < 5 мин бүртгэл & E2E: ~3 мин & $\checkmark$ \\
\hline
\end{tabular}
\end{table}

Системийн чанарын шаардлагууд туршилтаар бүрэн биелж, RQ1--RQ3 гурван асуултад тоон баталгаа олгосон.

\vspace{0.5cm}

%-------------------------------------------------------------------------------
%   SECTION - CONCLUSION
%-------------------------------------------------------------------------------

\section{Дүгнэлт}
\label{sec:conclusion}

Төгсөлтийн ажлын хүрээнд Монголын ЖДБ хоорондын захиалга, хүргэлтийн процессыг цахимжуулах \textbf{``Байгууллагын бүтээгдэхүүн хүргэлтийн нэгдсэн систем''}-ийг боловсруулж, туршилтын үр дүнгээр баталгаажуулав.

%===============================================================================
\subsection{Зорилтуудын хэрэгжилт}

Дэвшүүлсэн дөрвөн зорилтыг дараах байдлаар хэрэгжүүлсэн. Туршилтын үр дүнд үндэслэн зорилт бүрийн гүйцэтгэлд холбогдох нотолгоог нэмж оруулсан.

\begin{table}[H]
\centering
\renewcommand{\arraystretch}{1.3}
\caption{Зорилтуудын хэрэгжилтийн хураангуй}
\label{tab:objectives_summary}
\begin{tabular}{|c|p{3.5cm}|p{7.5cm}|c|}
\hline
\textbf{\#} & \textbf{Зорилт} & \textbf{Хэрэгжүүлэлт ба нотолгоо} & \textbf{Төлөв} \\
\hline
1 & Судалгаа, шинжилгээ & Uber Freight, Papa Logistics судалж, 4 архитектурын зөрүү тодорхойлсон & $\checkmark$ \\
\hline
2 & Архитектур, загварчлал & 11+ хүснэгт бүхий ӨС, 2 клиент (мобайл + веб) + Supabase backend, диаграммууд & $\checkmark$ \\
\hline
3 & Аюулгүй байдал & RLS бодлого, RBAC (4 түвшин), JWT validation, ePOD OTP; RLS penetration 6 тест болон OTP brute-force тестээр бүрэн нотлогдсон (Хүснэгт~\ref{tab:rls_pentest}, \S\ref{subsec:security_test}) & $\checkmark$ \\
\hline
4 & Функционал модулиуд & Захиалга, даалгаврын marketplace, realtime мэдэгдэл, каталог, ePOD, орлого, санхүүгийн тайлан; end-to-end урсгал баталгаажсан (\S\ref{subsec:functional_test}) & $\checkmark$ \\
\hline
\end{tabular}
\end{table}

%===============================================================================
\subsection{Харьцуулсан дүгнэлт}

Судалгаанд авч үзсэн хоёр платформтой харьцуулсан үр дүнг Хүснэгт~\ref{tab:unified_comparison}-д нэгтгэв. Туршилтаар баталгаажсан шинж чанарыг тухайн хүснэгтээр ишилсэн.

\begin{table}[H]
\centering
\renewcommand{\arraystretch}{1.3}
\caption{Гурван системийн нэгдсэн харьцуулалт}
\label{tab:unified_comparison}
\footnotesize
\begin{tabular}{|p{2.8cm}|p{2.8cm}|p{2.8cm}|p{3cm}|c|}
\hline
\textbf{Шалгуур} & \textbf{Uber Freight} & \textbf{Papa Logistics} & \textbf{Миний төсөл} & \textbf{Үнэлгээ} \\
\hline
Даалгавар хуваарилалт & ML алгоритм, recommendation & Оператор + санал & FCFS claim model, concurrency-safe (Хүснэгт~\ref{tab:concurrency_results}) & $\bigstar\bigstar$ \\
\hline
Үнэ тогтоолт & Динамик, ML-д суурилсан & Өрсөлдөөнт bid & Борлуулагч тогтоодог & $\bigstar\bigstar$ \\
\hline
Өгөгдлийн тусгаарлалт & Enterprise түвшин & Application-level & Database-level RLS, 6/6 pentest амжилттай (Хүснэгт~\ref{tab:rls_pentest}) & $\bigstar\bigstar\bigstar\bigstar$ \\
\hline
Audit trail & Execution log & Хязгаарлагдмал & Төлөв шилжилтийн timestamp; бүрэн audit log нь цаашдын ажил & $\bigstar\bigstar$ \\
\hline
API / гадаад integration & Бүрэн, TMS/ERP & Хязгаарлагдмал & RESTful суурь бүтэц & $\bigstar\bigstar\bigstar$ \\
\hline
Realtime хяналт & Location + tracking & Location tracking & Supabase Realtime мэдэгдэл & $\bigstar\bigstar\bigstar$ \\
\hline
Хүргэлтийн баталгаа & Дижитал signature & Тодорхойгүй & ePOD OTP код (brute-force туршигдсан, \S\ref{subsec:security_test}) & $\bigstar\bigstar\bigstar\bigstar$ \\
\hline
Performance & Enterprise түвшин & Үндсэн & p95 < 2с бүх endpoint-д (Хүснэгт~\ref{tab:endpoint_latency}) & $\bigstar\bigstar\bigstar\bigstar$ \\
\hline
Санхүүгийн тайлан & Enterprise түвшин & Үндсэн & Дэлгэрэнгүй (KPI, аналитик) & $\bigstar\bigstar\bigstar\bigstar$ \\
\hline
Өргөтгөх боломж & Cloud, horizontal scaling & Dispatcher bottleneck & Event-driven, horizontal & $\bigstar\bigstar\bigstar\bigstar$ \\
\hline
\end{tabular}
\end{table}

\textbf{Uber Freight-тэй харьцуулбал} чадамжийн ойролцоогоор 35--40\%-д хүрсэн. Uber нь олон жилийн ML, сая+ хэрэглэгчийн өгөгдөл, мянган инженертэй enterprise түвшний систем тул шууд харьцуулах боломжгүй боловч архитектурын суурь зарчим (autonomous claim-based хуваарилалт, database-level security, ePOD verification) ижил чиглэлтэй.

\textbf{Papa Logistics-тэй харьцуулбал} чадамжийн ойролцоогоор 70--75\%-д хүрсэн. Архитектурын хувьд (RLS, audit timestamp, операторгүй marketplace) давуу; харин бодит зах зээлийн туршлага, төлбөрийн integration, location tracking-оор дутмаг.

%===============================================================================
\subsection{Хийж чадаагүй зүйлс ба цаашдын ажил}

Хийгдсэн ажлын дэлгэрэнгүйг §\ref{subsec:results_summary}-ийн Хүснэгт~\ref{tab:objectives_summary} (Зорилтуудын хэрэгжилт) болон Хүснэгт~\ref{tab:requirement_verification} (Шаардлагын биелэлт)-д нэгтгэсэн. Энэ хэсэгт одоогийн хувилбарт хамрагдаагүй, цаашид хөгжүүлэх чиглэлүүдийг тусгайлан дурдав.

\textbf{Захиалагчийн (customer) талын апп.} Одоогийн хувилбарт борлуулагчийн веб апп, хүргэгчийн мобайл апп гэсэн хоёр клиент хөгжүүлсэн. Эцсийн хэрэглэгч (бараа худалдан авагч) нь платформын two-sided marketplace-ийн эрэлтийн тал боловч уг ажлын хүрээнд тусдаа апп хэрэгжээгүй. Одоогоор захиалагч нь борлуулагчийн байгууллагаар дамжуулан захиалга өгдөг хэлбэрээр загварчилсан (\texttt{orders} хүснэгтэд \texttt{customer\_id} талбар бэлтгэгдсэн). Цаашид захиалагчийн мобайл апп (бүтээгдэхүүн хайх, сагс, төлбөр, хүргэлт хянах) нэмэх нь зайлшгүй.

\textbf{Цаашид хөгжүүлэх бусад чиглэл:}
\begin{itemize}
    \item Бодит төлбөрийн integration (QPay, SocialPay, банкны API)
    \item Realtime GPS location tracking
    \item ML-д суурилсан даалгаврын recommendation
    \item Route optimization алгоритм
    \item Push notification систем
    \item Бодит хэрэглэгчтэй pilot туршилт
    \item 100+ concurrent user-тэй өргөтгөсөн load test (read replica, per-tenant sharding нөхцөлд)
\end{itemize}

%===============================================================================
\subsection{Төгсгөл}

Тус төсөл нь Монгол Улсын B2B хүргэлтийн зах зээлд operator-гүй, autonomous FCFS claim хуваарилалт болон database-level security-д суурилсан шинэ архитектурын шийдлийг concurrency, penetration, load туршилтаар тоон үзүүлэлтээр баталгаажуулсан. Цаашид §\ref{subsec:results_summary}-д дурдсан чадамжуудыг нэмж хөгжүүлснээр production-д нэвтрэх боломжтой систем болно.



%===============================================================================
% Ном зүй
%===============================================================================
\newpage
\phantomsection
\addcontentsline{toc}{section}{Ном зүй}
\printbibliography[title={Ном зүй}]

%===============================================================================
% Товчлол (Summary) — two-column condensed thesis summary at the very end
%===============================================================================
%-------------------------------------------------------------------------------
%   SUMMARY (ТОВЧЛОЛ) — Two-column condensed thesis summary
%-------------------------------------------------------------------------------

\newpage
\phantomsection
\addchaptertocentry{Товчлол}
\thispagestyle{plain}

\vspace*{-6mm}
\begin{center}
{\headingfont\large \ttitle}\\[1mm]
\end{center}

\begin{flushright}
\deptname\\
Оюутан: \authorcode, \authorshort\\
Удирдагч багш: \supervisor
\end{flushright}

\vspace{1mm}
\setcounter{table}{0}
\setcounter{figure}{0}
\renewcommand{\thetable}{\arabic{table}}
\renewcommand{\thefigure}{\arabic{figure}}
\setlength{\columnsep}{7mm}
\begin{multicols}{2}
\setlength{\parskip}{5pt}
\sloppy

\noindent{\headingfont 1. Удиртгал}

Монгол Улсын жижиг, дунд бизнесийн байгууллагууд (ЖДБ) хоорондоо бараа захиалах, хүргэх үйл явцаа одоо хүртэл голдуу утас, мессеж, цаасан бичгээр гүйцэтгэсээр байна. Одоогийн практикт (i) жолооч бүртэй тусад нь утсаар ярьж захиалга бүрт 10–15 минут зарцуулдаг, (ii) хүргэлтийн баримтыг гар бичмэлээр хөтлөх нь андуурал, маргааныг нэмэгдүүлдэг, (iii) байгууллага бүр өөрийн цалинжуулдаг хүргэлтийн багтай байхыг шаарддаг зэрэг асуудал байсаар байна. Ийм гараар зохицуулдаг хэлбэр нь цаг алдагдуулж, зохион байгуулалтын алдаа үүсгэж, явцыг хянахад төвөгтэй болгодог. Монголын зах зээлд B2B (business-to-business буюу байгууллага хоорондын) хүргэлтийн нэгдсэн шийдэл хараахан төлөвшөөгүй бөгөөд одоогийн тоглогчид голдуу dispatcher (захиалга хуваарилагч ажилтан)-т тулгуурласан эсвэл bid-based (үнийн саналын уралдаан) нэг талтай загвартай хэвээр байна.

\textbf{Operator-гүй загвар.} ``Operator-гүй'' (оператор-free) гэдэг нь систем дотор захиалгыг гараар хуваарилдаг төвлөрсөн ажилтан (диспетчер) байхгүй, харин хүргэгчид (courier) боломжит даалгаврын жагсаалтаас өөрсдөө сонгон авдаг (self-service claim) загварыг хэлнэ. Энэ нь хүний хүчин зүйлийн алдаа, bottleneck-ийг (хүлцэгдэх хүзүү) арилгана.

\textbf{Системийн оролцогчид.} Платформ дөрвөн үндсэн оролцогчтой: (1) \emph{Байгууллага} (organization/tenant) — бүтээгдэхүүн бүртгэж, захиалга хүлээн авч, хүргэлт зарладаг; (2) \emph{Хэрэглэгч} (customer) — захиалга өгч, ePOD OTP-ээр хүлээн авсан баталгаажуулалт хийдэг; (3) \emph{Хүргэгч} (courier) — даалгаврыг claim (барьж авах) хийж, хүргэлтийг гүйцэтгэдэг; (4) \emph{Админ} (admin) — байгууллагын бүртгэл болон системийн ерөнхий тохиргоог хянадаг.

Судалгааны зорилго нь:
\begin{itemize}
    \item Operator-гүй (төвлөрсөн зохицуулагчгүй), multi-tenant (олон байгууллагыг нэг системд тусгаарлан үйлчлэх) B2B хүргэлтийн платформыг Supabase (PostgreSQL + PostgREST + RLS буюу мөрийн түвшний аюулгүй байдал) технологи дээр боловсруулах,
    \item Concurrency (зэрэгцээ хандалт) хамгаалалтыг database (өгөгдлийн сан) түвшинд шийдэж баталгаажуулах,
    \item Multi-tenant data isolation-ыг (байгууллага хоорондын өгөгдлийн тусгаарлалтыг) нотлох,
    \item ePOD (electronic Proof of Delivery буюу цахим хүргэлтийн нотолгоо) OTP (One-Time Password, нэг удаагийн нууц үг) верификацыг аюулгүй болгох,
    \item Туршилтаар системийн performance-ийг (ажиллагааны хурд) үнэлэх.
\end{itemize}

Судалгааны асуултууд (RQ --- Research Question) дараах гурван чиглэлд төвлөрсөн: RQ1 — FCFS (First-Come-First-Served, эхэлж ирсэн нь эхэлж үйлчлүүлэх) claim загварын concurrency correctness (зэрэгцээ хандалтын зөв эсэх); RQ2 — RLS-ийн multi-tenant isolation найдвартай байдал; RQ3 — BaaS-mediated (Backend-as-a-Service буюу backend-ийг үйлчилгээ байдлаар хүргэх платформоор зуучилсан) архитектурын production viability (үйлдвэрлэлд ашиглахад тэнцэх байдал). Эдгээрийг concurrency test (зэрэгцээ хандалтын сорил), load test (ачааллын сорил), security penetration test (аюулгүй байдлын нэвтрэлтийн сорил), functional E2E (end-to-end буюу эхнээс эцэс хүртэл) test гэсэн 4 төрлийн туршилтаар хэмжив.

\noindent{\headingfont 2. Судалгааны сэдвийн онол, өнөөгийн байдал}

B2B хүргэлтийн платформын онолын суурь нь таван үндсэн чиглэлд тулгуурласан: last-mile (сүүлчийн миль буюу эцсийн хүргэлтийн үе) логистик, two-sided marketplace (хоёр талт зах зээл), multi-tenant isolation, concurrency control (зэрэгцээ хандалтын хяналт), BaaS (Backend-as-a-Service) архитектур.

Last-mile хүргэлт нь нийт тээврийн зардлын 13--75\%-ийг эзэлдэг тул crowdsourced (олон нийтийн оролцоотой) загвар нь уламжлалт өөрийн парктай системийг орлох шинэ чиг хандлага болсон. Two-sided marketplace онолоор (Rochet \& Tirole 2003) B2B платформын гол саад нь liquidity (хоёр талын оролцогчдын хүрэлцэх нягтрал) тул cold-start (хэрэглэгч цөөн эхний үе) шатанд FCFS загвар давуу талтай.

Multi-tenant өгөгдлийн тусгаарлалтыг Chong (2006), Bezemer \& Zaidman (2010) нар database түвшинд хэрэгжүүлэхийг зөвлөсөн. Үүнийг PostgreSQL-ийн Row Level Security (RLS, мөрийн түвшний хандалтын хяналт) нь хүснэгтийн мөр бүрд \texttt{USING}/\texttt{WITH CHECK} policy (бодлого) хавсарган хангана.

Concurrency control-ыг Bernstein (1987), Kleppmann (2017) нар pessimistic (урьдчилан түгжих) / optimistic (сүүлд шалгах) хоёр хандлагад ангилсан. Write-heavy (бичих үйлдэл давамгайлсан) workload-д pessimistic илүү найдвартай тул PostgreSQL-ийн \texttt{SELECT ... FOR UPDATE SKIP LOCKED} (түгжигдсэн мөрийг алгасах SQL хэлбэр) нь race condition-ыг (зэрэгцээ мөргөлдөөнийг) 0 магадлалаар шийднэ.

Харьцуулсан судалгаанд Uber Freight (ML буюу машин сургалтад суурилсан enterprise (томоохон корпорацын) платформ) болон Papa Logistics (bid-based буюу илгээгч захиалга нийтлээд жолооч нар үнийн санал илгээж, илгээгч харьцуулан сонгодог Монголын платформ) хоёрыг авч үзсэн. Эдгээрээс dispatcher bottleneck (захиалга хуваарилагчид төвлөрсөн хүлцэгдэх хүзүү), data isolation сул, audit trail (үйлдлийн түүхийн бүртгэл) хязгаарлагдмал, API (Application Programming Interface, програмчлалын интерфэйс) integration сул гэсэн 4 архитектурын зөрүүг тодорхойлсон.

\noindent{\headingfont 3. Судалгааны арга зүй}

Судалгааг Design Science Research (загварчлалын шинжлэх ухааны судалгаа) арга барилаар 3 үе шатаар (загварчлал → хөгжүүлэлт → туршилт) Agile (уян хатан, богино давталттай хөгжүүлэлтийн арга) давталттай хэрэгжүүлэв. Платформ нь хоёр бие даасан клиент апп (Хүргэгчийн React Native мобайл апп, Борлуулагчийн Next.js веб апп) болон тэдгээртэй HTTPS/WebSocket-ээр холбогдсон нэгдсэн Supabase backend-ээс бүрдэнэ (Зураг~\ref{fig:tovchlol_overview}).

\begin{figure}[H]
\centering
\resizebox{\columnwidth}{!}{%
\begin{tikzpicture}[
  font=\scriptsize,
  appbox/.style={rectangle, rounded corners=4pt, draw=black!55, thick,
    minimum width=4.2cm, minimum height=2.9cm, align=left, inner sep=5pt, text width=3.8cm},
  courierbox/.style={appbox, fill=green!15, draw=green!55!black},
  merchantbox/.style={appbox, fill=blue!12, draw=blue!55!black},
  backendbox/.style={rectangle, rounded corners=4pt, draw=black!60, thick, fill=gray!15,
    minimum width=9.2cm, minimum height=1.8cm, align=center, inner sep=5pt},
  subbox/.style={rectangle, rounded corners=3pt, draw=blue!70!black, thick, fill=blue!20,
    minimum width=2.5cm, minimum height=0.55cm, align=center, font=\tiny\bfseries, inner sep=3pt},
  arrow/.style={<->, >=Latex, thick, draw=gray!65},
  arrlabel/.style={fill=white, inner sep=1pt, font=\tiny},
]
\node[courierbox] (courier) at (-2.7, 2.3) {%
  {\footnotesize\bfseries Хүргэгчийн мобайл апп}\\[1pt]
  {\tiny\itshape\textcolor{black!65}{React Native / Expo}}\\[2pt]
  $\bullet$~Бүртгэл, баталгаажуулалт\\
  $\bullet$~Нээлттэй даалгавар\\
  $\bullet$~Даалгавар claim\\
  $\bullet$~Хүргэлт гүйцэтгэх\\
  $\bullet$~KYC илгээх};
\node[merchantbox] (merchant) at (2.7, 2.3) {%
  {\footnotesize\bfseries Борлуулагчийн веб апп}\\[1pt]
  {\tiny\itshape\textcolor{black!65}{Next.js}}\\[2pt]
  $\bullet$~Бүтээгдэхүүн удирдах\\
  $\bullet$~Даалгавар үүсгэх\\
  $\bullet$~Зах зээлд нийтлэх\\
  $\bullet$~Гүйцэтгэл хянах\\
  $\bullet$~Тайлан, аналитик};
\node[backendbox] (backend) at (0, -1.35) {};
\node[font=\footnotesize\bfseries] at (0, -0.75) {Backend үйлчилгээ (Supabase)};
\node[font=\tiny\itshape, text=black!65] at (0, -1.1) {Нийтлэг API, өгөгдлийн сан, бодит цагийн холболт};
\node[subbox] at (-3.1, -1.65) {RESTful API};
\node[subbox] at (0, -1.65) {PostgreSQL + RLS};
\node[subbox] at (3.1, -1.65) {WebSocket};
\draw[arrow] (courier.south) -- node[arrlabel, pos=0.55] {API / WS} ([xshift=-2.7cm]backend.north);
\draw[arrow] (merchant.south) -- node[arrlabel, pos=0.55] {API / WS} ([xshift=2.7cm]backend.north);
\end{tikzpicture}
}
\caption{Системийн ерөнхий бүтэц: 2 клиент апп + нэгдсэн backend}
\label{fig:tovchlol_overview}
\end{figure}

Архитектурын гол онцлог нь \textbf{BaaS-mediated architecture} (BaaS-ээр зуучилсан архитектур) бөгөөд уламжлалт three-tier (гурван давхарга бүхий сонгодог client–server–database) биш, Supabase BaaS платформоор бүх backend (серверийн тал) үйлчилгээг зуучилсан гурван давхаргат бүтэц юм (Зураг~\ref{fig:tovchlol_arch}).

\begin{figure}[H]
\centering
\resizebox{\columnwidth}{!}{%
\begin{tikzpicture}[
  layer/.style={draw, rounded corners, minimum width=10cm, align=center, font=\footnotesize},
  arrow/.style={->, >=stealth, thick},
]
\node[layer, fill=blue!10, minimum height=1.1cm] (pres) at (0, 0) {
  \textbf{Client layer}\\
  React Native (мобайл) $|$ Next.js (веб) $|$ SDK
};
\node[layer, fill=orange!10, minimum height=1.6cm] (baas) at (0, -2.0) {
  \textbf{BaaS layer --- Supabase}\\
  PostgREST $|$ GoTrue (JWT) $|$ Realtime\\
  Storage $|$ Edge Functions $|$ RPC
};
\node[layer, fill=green!10, minimum height=1.1cm] (data) at (0, -3.9) {
  \textbf{Data layer}\\
  PostgreSQL $|$ RLS $|$ Triggers $|$ SECURITY DEFINER
};
\draw[arrow] (pres.south) -- (baas.north) node[midway, right, font=\tiny] {HTTPS / WS};
\draw[arrow] (baas.south) -- (data.north) node[midway, right, font=\tiny] {SQL / PL/pgSQL};
\end{tikzpicture}
}
\caption{BaaS-mediated архитектурын гурван давхаргат бүтэц}
\label{fig:tovchlol_arch}
\end{figure}

Client layer (клиент/хэрэглэгчийн давхарга) нь хэрэглэгчийн интерфэйсийг хариуцна. BaaS layer нь Supabase платформоор API-г автоматаар үүсгэж, JWT (JSON Web Token --- нууцлан гарын үсэг зурсан хандалтын токен) authentication (баталгаажуулалт), realtime (шууд дамжих) WebSocket, storage (файл хадгалалт) болон RPC (Remote Procedure Call --- алсын функц дуудлага) дуудлагыг нэгтгэнэ. Data layer нь PostgreSQL + RLS + триггер (автомат ажилладаг database функц) + \texttt{SECURITY DEFINER} (өндөр эрхээр гүйцэтгэгддэг) функцээр бизнесийн логикийг database дотор хэрэгжүүлнэ.

Өгөгдлийн түвшинд 11 хүснэгт бүхий schema (өгөгдлийн бүтэц) зохион байгуулж, multi-tenant isolation-ыг \texttt{org\_id} талбар дээр суурилсан RLS-ээр, зэрэгцээ claim-ийг \texttt{FOR UPDATE SKIP LOCKED} lock-оор (түгжих механизм) хангасан. Бизнесийн логикийн гол функцүүд: \texttt{claim\_delivery\_task} (атомик --- хуваагдашгүй байдлаар даалгавар оноох), \texttt{verify\_epod\_otp} (bcrypt --- аюулгүй нууц үг хэш хийх алгоритмаар OTP баталгаажуулах), мөн автомат төлөв шилжилт, орлого бүртгэлийн триггерүүд.

Proof of Delivery (ePOD --- хүргэлтийн баталгаа)-г 6 оронтой OTP код (bcrypt hash --- нууцалсан хэш, 10 мин expiry буюу хүчинтэй хугацаа, 5 оролдлогын хязгаар) хэлбэрээр хэрэгжүүлсэн. Client-server contract-ыг (клиент–сервер хоорондын өгөгдлийн гэрээ) хамгаалах schema drift prevention (схемийн зөрчил үүсэхээс сэргийлэх) стратегийг Supabase CLI (командын мөрийн хэрэгсэл)-ийн auto-generated (автоматаар үүсгэсэн) TypeScript төрөл + CI (Continuous Integration --- тасралтгүй нэгдэл) шатны харьцуулалтаар хангасан.

\textbf{Системийн үндсэн урсгал.} Бүртгэлээс орлого бүртгэл хүртэлх үйл явцыг 8 алхамт дараалалд зохион байгуулсан (Зураг~\ref{fig:tovchlol_flow}). Энэ урсгал нь веб апп (1–4 алхам), мобайл апп (5–7 алхам), database триггер (8 алхам) гэсэн 3 бүрэлдэхүүнийг нэг гинжид уялдуулсан.

\textbf{Өгөгдлийн урсгалын контекст (DFD Level 0).} Платформ нь 4 гадаад оролцогчтой харилцдаг (Зураг~\ref{fig:tovchlol_dfd}): борлуулагч бүтээгдэхүүн/даалгавар илгээж, гүйцэтгэл–орлогын тайлан хүлээж авна; хүргэгч нээлттэй даалгаврыг харж, KYC ба ePOD илгээнэ; имэйл үйлчилгээ OTP/мэдэгдэл дамжуулна; админ KYC-г батлан, бодлого тохируулна.

\begin{figure}[H]
\centering
\resizebox{0.85\columnwidth}{!}{%
\begin{tikzpicture}[
  font=\footnotesize,
  ent/.style={draw, rectangle, rounded corners=2pt, fill=blue!5,
    minimum width=1.6cm, minimum height=0.7cm, align=center, inner sep=2pt},
  core/.style={draw, ellipse, fill=green!20,
    minimum width=2.8cm, minimum height=1.2cm, align=center, inner sep=2pt},
  lbl/.style={font=\scriptsize, inner sep=1pt},
  flow/.style={->, >=stealth, semithick}
]
% Nodes
\node[core] (p) at (0,0) {Хүргэлтийн\\платформ};
\node[ent] (mail)    at (0, 2.0)  {Имэйл\\сервис};
\node[ent] (seller)  at (-3.8, 0) {Борлуулагч};
\node[ent] (courier) at (3.8, 0)  {Хүргэгч};
\node[ent] (admin)   at (0, -2.0) {Админ};

% Mail (vertical, top)
\draw[flow] ($(p.north)+(0.2,0)$)  -- node[lbl, right]{OTP, мэдэгдэл} ($(mail.south)+(0.2,0)$);
\draw[flow] ($(mail.south)+(-0.2,0)$) -- node[lbl, left]{баталгаа} ($(p.north)+(-0.2,0)$);

% Seller (horizontal, left)
\draw[flow] (seller.east |- 0,0.15) -- node[lbl, above]{даалгавар} (p.west |- 0,0.15);
\draw[flow] (p.west |- 0,-0.15) -- node[lbl, below]{тайлан} (seller.east |- 0,-0.15);

% Courier (horizontal, right)
\draw[flow] (p.east |- 0,0.15) -- node[lbl, above]{даалгавар, OTP} (courier.west |- 0,0.15);
\draw[flow] (courier.west |- 0,-0.15) -- node[lbl, below]{KYC, ePOD} (p.east |- 0,-0.15);

% Admin (vertical, bottom)
\draw[flow] ($(p.south)+(-0.2,0)$) -- node[lbl, left]{KYC хүсэлт} ($(admin.north)+(-0.2,0)$);
\draw[flow] ($(admin.north)+(0.2,0)$) -- node[lbl, right]{батлах, бодлого} ($(p.south)+(0.2,0)$);
\end{tikzpicture}
}
\caption{DFD Level 0 --- платформын контекст}
\label{fig:tovchlol_dfd}
\end{figure}

\begin{figure}[H]
\centering
\resizebox{0.95\columnwidth}{!}{%
\begin{tikzpicture}[
  font=\scriptsize,
  stepbox/.style={rectangle, rounded corners=3pt, draw=green!60!black, fill=green!25,
    line width=0.7pt, minimum width=2.5cm, minimum height=0.7cm, align=center, inner sep=2pt},
  flowarrow/.style={->, >=stealth, thick, draw=green!60!black}
]
  \node[stepbox] (s1) at (0, 0)     {\textbf{1.} Байгууллага\\бүртгүүлэх};
  \node[stepbox] (s2) at (3.0, 0)   {\textbf{2.} Бүтээгдэхүүн\\нэмэх};
  \node[stepbox] (s3) at (3.0, -1.2){\textbf{3.} Захиалга\\хүлээн авах};
  \node[stepbox] (s4) at (0, -1.2)  {\textbf{4.} Даалгавар\\нийтлэх};
  \node[stepbox] (s5) at (0, -2.4)  {\textbf{5.} Даалгавар\\claim};
  \node[stepbox] (s6) at (3.0, -2.4){\textbf{6.} Хүргэлт\\гүйцэтгэх};
  \node[stepbox] (s7) at (3.0, -3.6){\textbf{7.} ePOD OTP\\баталгаажуулалт};
  \node[stepbox] (s8) at (0, -3.6)  {\textbf{8.} Орлого\\бүртгэх};
  \draw[flowarrow] (s1) -- (s2);
  \draw[flowarrow] (s2) -- (s3);
  \draw[flowarrow] (s3) -- (s4);
  \draw[flowarrow] (s4) -- (s5);
  \draw[flowarrow] (s5) -- (s6);
  \draw[flowarrow] (s6) -- (s7);
  \draw[flowarrow] (s7) -- (s8);
\end{tikzpicture}
}
\caption{Системийн үндсэн урсгал: 8 алхамт үе шат}
\label{fig:tovchlol_flow}
\end{figure}

\textbf{Аюулгүй байдлын олон давхаргат хамгаалалт.} Системд 5 түвшний security (аюулгүй байдлын) давхарга хэрэгжсэн: (1) HTTPS/TLS (сүлжээн дэх шифрлэгдсэн холболтын протокол) сүлжээний шифрлэлт, (2) Supabase Auth + JWT, (3) Next.js middleware (хүсэлт хүлээн авах завсрын давхарга) / Navigation guard (хуудсын хандалтыг шалгагч), (4) database-level RLS бодлого (\texttt{org\_id}, \texttt{courier\_id}, \texttt{auth.uid()}), (5) \texttt{SECURITY DEFINER} RPC функц дэх бизнесийн дүрмийн шалгалт.

\begin{table}[H]
\centering
\caption{Технологийн стекийн хураангуй}
\label{tab:tovchlol_stack}
\scriptsize
\setlength{\tabcolsep}{2.5pt}
\begin{tabular}{|l|l|}
\hline
\textbf{Давхарга} & \textbf{Технологи} \\ \hline
Мобайл апп & React Native + Expo 54, TS 5.9 \\ \hline
Веб апп & Next.js 16.1, React 19.2 \\ \hline
Backend & Supabase (PostgreSQL 15+) \\ \hline
Security & RLS, JWT, bcrypt ePOD OTP \\ \hline
\end{tabular}
\end{table}

Туршилтын орчин: Supabase Pro tier (төлбөртэй үндсэн багц), PostgreSQL 15.6, PgBouncer transaction pooling (холболтын нөөцлөлт). Өгөгдлийн синтетик (хиймэл үүсгэсэн) багц: 20 байгууллага, 200 хэрэглэгч, 1,000 захиалга, 1,000 даалгавар. Ашигласан хэрэгсэл: \texttt{@supabase/supabase-js} (JavaScript client сан), \texttt{k6 0.54} (ачааллын сорилын хэрэгсэл), \texttt{psql} + \texttt{SET request.jwt.claims} (JWT claim-ийг шууд тохируулах SQL команд).

\noindent{\headingfont 4. Судалгааны үр дүн, нэгтгэсэн дүгнэлт}

\textbf{Concurrency test.} $N \in \{10, 50, 100\}$ хүргэгч нэг даалгаврыг зэрэг claim (барьж авах) хийх 30 сорилын туршид \textbf{давхардсан claim нэг ч удаа гараагүй} (100\% амжилт). Дундаж хариу хугацаа $N{=}10$-д 124~ms, $N{=}50$-д 287~ms, $N{=}100$-д 512~ms болсон нь O($n$) шугаман хэмжээтэй, deadlock (харилцан түгжрэл)-гүй pessimistic locking-ийн онолын хүлээгдсэн зан төлөвтэй нийцнэ.

\textbf{Load test.} k6-аар 100 VU (Virtual User --- виртуал хэрэглэгч), 30 мин ачаалал өгөхөд бүх endpoint-ийн (API холбох цэг) p95 latency (95\%-ийн хариу хугацаа) 2 секундын зорилтот хязгаарт багтсан:
\end{multicols}

\begin{figure}[H]
\centering
\resizebox{0.85\columnwidth}{!}{%
\begin{tikzpicture}
\begin{axis}[
    ybar,
    bar width=11pt,
    width=0.92\textwidth,
    height=6.2cm,
    enlarge x limits=0.12,
    ymin=0, ymax=900,
    ylabel={Хариу хугацаа (ms)},
    symbolic x coords={available\_tasks,claim\_task,verify\_otp,orders,dashboard},
    xtick=data,
    x tick label style={font=\scriptsize, rotate=15, anchor=east},
    ymajorgrids=true,
    grid style={gray!30},
    legend style={at={(0.5,-0.32)}, anchor=north, legend columns=3, font=\footnotesize},
    ylabel style={font=\small},
    tick label style={font=\footnotesize}
]
\addplot[fill=blue!40, draw=blue!60!black] coordinates {
    (available\_tasks,145) (claim\_task,287) (verify\_otp,198) (orders,167) (dashboard,234)
};
\addplot[fill=orange!60, draw=orange!70!black] coordinates {
    (available\_tasks,312) (claim\_task,621) (verify\_otp,287) (orders,389) (dashboard,567)
};
\addplot[fill=red!50, draw=red!70!black] coordinates {
    (available\_tasks,456) (claim\_task,834) (verify\_otp,345) (orders,512) (dashboard,789)
};
\legend{p50, p95, p99}
\end{axis}
\end{tikzpicture}
}
\caption{k6 load test — endpoint latency хуваарилалт (100 VU, 30 мин)}
\label{fig:tovchlol_latency}
\end{figure}

\begin{multicols*}{2}

\textbf{Security penetration test.} Org~A хэрэглэгчийн JWT-ээр Org~B-гийн өгөгдөлд хандах 6 хувилбар (cross-tenant буюу байгууллага хооронд SELECT/UPDATE, SQL injection --- хорлонтой SQL оруулах, pending courier access --- зөвшөөрөгдөөгүй хүргэгчийн хандалт, JWT forgery --- токен хуурамчлах) бүгд амжилтгүй болсон (6/6). OTP brute-force (дараалан оролдох халдлага) туршилтад 5 дахь алдаатай оролдлогын дараа сесс блоклогдож, brute-force амжилтын онолын магадлал $5 \times 10^{-4}\%$ болж буурсан.

\textbf{Functional E2E.} Байгууллагын бүртгэл → бүтээгдэхүүн нэмэх → захиалга → даалгавар нийтлэх → хүргэгч claim → pickup (авалт) → delivery (хүргэлт) → ePOD OTP → орлого бүртгэл гэсэн 8 алхамт урсгал зөв дарааллаар биелэв. Timestamp invariant (цагийн тэмдгийн өөрчлөгдөөгүй дараалал --- \texttt{assigned\_at} < \texttt{picked\_up\_at} < \texttt{delivered\_at} < \texttt{completed\_at}) зөрөлгүй, хүргэлтийн орлого автомат бүртгэгдсэн.

\begin{table}[H]
\centering
\caption{Харьцуулсан үнэлгээний хураангуй}
\label{tab:tovchlol_comparison}
\scriptsize
\begin{tabular}{|l|c|}
\hline
\textbf{Платформ} & \textbf{Чадамжийн хувь} \\ \hline
Uber Freight & $\sim$35--40\% \\ \hline
Papa Logistics & $\sim$70--75\% \\ \hline
\end{tabular}
\end{table}

\textbf{Хязгаарлалт.} Туршилт нь Supabase Pro tier-ийн нөөцийн хязгаарт, 100 VU хүртэлх concurrent (зэрэгцээ) ачаалалд, synthetic (хиймэл үүсгэсэн) өгөгдөл дээр, single-region (нэг бүсийн сервер) тохиргоонд явагдсан. 500+ VU, multi-region (олон бүсийн), 3G/4G сүлжээний нөхцөл үнэлэгдээгүй. Бодит GPS tracking (байршил хяналт), QPay/SocialPay төлбөрийн integration (холболт), ML recommendation (машин сургалтаар санал болгох систем) нь future work (ирээдүйн судалгаа)-т үлдсэн.

\textbf{Ерөнхий дүгнэлт.} Судалгааны гол дүгнэлт нь B2B хүргэлтийн платформын найдвартай байдлыг database түвшинд хэрэгжүүлсэн concurrency control (\texttt{FOR UPDATE SKIP LOCKED}) ба RLS гэсэн хоёр тулгуурт механизмд суурилан бүтээх боломжтойг туршилтын тоогоор нотолсон явдал юм. FCFS claim загвар, row-level isolation, bcrypt-д суурилсан ePOD OTP гэсэн гурван архитектурын шийдэл хамтдаа ажилласнаар multi-tenant орчинд race condition, cross-tenant access, OTP brute-force гэсэн гурван гол эрсдэлийг бодитоор шийдсэн.

\noindent{\headingfont 5. Ном зүй}

\small
Судалгааны гол эх сурвалжууд:
\begin{enumerate}
    \itemsep=0pt
    \item Chong, F., Carraro, G., \& Wolter, R. (2006). Multi-tenant data architecture.
    \item Bezemer, C.-P., \& Zaidman, A. (2010). Multi-tenant SaaS applications: Maintenance dream or nightmare?
    \item Ferraiolo, D. F., \& Kuhn, D. R. (1992). Role-based access controls.
    \item Bernstein, P. A., Hadzilacos, V., \& Goodman, N. (1987). Concurrency Control and Recovery in Database Systems.
    \item Kleppmann, M. (2017). Designing Data-Intensive Applications. O'Reilly.
    \item Provos, N., \& Mazi\`eres, D. (1999). A future-adaptable password scheme (bcrypt).
    \item Jones, M., Bradley, J., \& Sakimura, N. (2015). JSON Web Token (RFC 7519).
    \item Baldini, I., et al. (2017). Serverless computing: Current trends and open problems.
    \item Rochet, J.-C., \& Tirole, J. (2003). Platform competition in two-sided markets.
    \item Alnaggar, A., et al. (2021). Crowdsourced delivery: A review.
\end{enumerate}
\normalsize

\end{multicols*}


\end{document}